\documentclass[a4paper,onecolumn,11pt]{quantumarticle}
\pdfoutput=1
\usepackage[utf8]{inputenc}
\usepackage[english]{babel}
\usepackage[T1]{fontenc}
\usepackage{xcolor}
\usepackage[compatibility=false]{caption}
\usepackage{mathdots}
\usepackage{amsmath}
\usepackage{comment}
\usepackage{mathtools}
\usepackage{amssymb}
\usepackage{amsfonts}
\usepackage{amsthm}
\usepackage{mathrsfs}
\usepackage{footnote}
\usepackage{graphicx} 
\usepackage{stmaryrd}
\usepackage{dcolumn}
\usepackage{tikz-cd}
\usepackage[colorlinks=true
,urlcolor=teal
,anchorcolor=blue
,citecolor=teal
,filecolor=blue
,linkcolor=teal
,menucolor=blue
,linktocpage=true
,pdfproducer=medialab
,pdfa=true
]{hyperref}


\newtheorem{thm}{Theorem}
\newtheorem{lem}{Lemma}
\newtheorem{conj}{Conjecture}
\newtheorem{cor}{Corollary}
\newtheorem{defn}{Definition}
\newtheorem{prop}{Proposition}
\newtheorem{cl}{Claim}
\newtheorem{rem}{Remark}

% Define a generic theorem-like structure
\newtheorem{innercustomgeneric}{\customgenericname}
\providecommand{\customgenericname}{}

% Define the specific environment for propositions with custom numbering
\newenvironment{mprop}[1]{%
  \renewcommand\customgenericname{Proposition}%
  \renewcommand\theinnercustomgeneric{#1}%
  \innercustomgeneric
}{\endinnercustomgeneric}

\newenvironment{mthm}[1]{%
  \renewcommand\customgenericname{Theorem}%
  \renewcommand\theinnercustomgeneric{#1}%
  \innercustomgeneric
}{\endinnercustomgeneric}


\newenvironment{mlemma}[1]{%
  \renewcommand\customgenericname{Lemma}%
  \renewcommand\theinnercustomgeneric{#1}%
  \innercustomgeneric
}{\endinnercustomgeneric}

\newcommand{\dd}{\mathrm{d}}
\newcommand{\ket}[1]{| #1 \rangle}
\newcommand{\kket}[1]{| #1 \rangle\rangle}
\newcommand{\bra}[1]{\langle #1 |}
\newcommand{\braket}[2]{\left\langle #1 \mid #2 \right\rangle}
\newcommand{\ketbra}[2]{\left|#1\right\rangle\!\!\left\langle #2\right|}
\def\N{\mathcal{N}}
\def\eps{\varepsilon}
\def\O{\mathcal{O}}
\def\h{\mathcal{H}}
\def\k{\mathcal{K}}
\def\id{\mathbb{I}}
\newcommand*{\Tr}{\mathrm{Tr}}
\def\tr{\mathrm{tr}}
\def\D{\mathcal{D}}
\def\E{\mathcal{E}}
\def\A{\mathcal{A}}
\def\C{\mathcal{C}}
\def\map{\mathcal}
\def\set{\mathsf}

\newcommand{\JW}[1]{\textcolor{teal}{[JW:\ #1]}}

\newcommand{\YY}[1]{\textcolor{blue}{[YY:\ #1]}}


\begin{document}

\title{Quantum minimum description of density matrices}

\author{Patrick Hayden}
\affiliation{Leinweber Institute for Theoretical Physics, Stanford University, Stanford, CA 94305, USA}

\author{Alexander Maloney}
\affiliation{Department of Physics, Syracuse University, Syracuse, NY 13244}
\affiliation{Institute for Quantum and Information Sciences, Syracuse University, Syracuse, NY 13244, USA}

\author{Jinzhao Wang}
\email{wang.jinzhao226@gmail.com}
\affiliation{Leinweber Institute for Theoretical Physics, Stanford University, Stanford, CA 94305, USA}

\author{Yuxiang Yang}
\affiliation{QICI Quantum Information and Computation Initiative, School of Computing and Data Science, The University of Hong Kong, Pokfulam Road, Hong Kong, China}
 
\maketitle

\begin{abstract}
We study the problem of finding the quantum minimal description of many copies of a density matrix $\rho^{\otimes n}$. Operationally, this corresponds to compressing $\rho^{\otimes n}$ without preserving its purification. We determine the optimal compression scheme with the optimal rate at the second order and establish its connection to free entropy~\cite{paper1}. 
Our main technical tool is a generalization of Werner’s cloning map to arbitrary $\mathrm{GL}(d,\mathbb C)$ irreducible representations. We introduce this generalized cloning map and analyze its structural properties. In particular, we show that it approximately maps the irrep state $\rho_\mu$ to $\rho_\nu$ with high fidelity whenever $\mu$ and $\nu$ are close, which is the key ingredient enabling our compression scheme.
\end{abstract}


\section{Introduction}\label{sec:intro}

%Motivation
Schumacher compression is a foundational result in quantum information theory~\cite{schumacher1995quantum}, and incidentally it is also where the name ``qubit'' was coined. In this paper, we study a variant of Schumacher compression, and through that we confirm its connection with a new entropy measure as claimed in a companion paper~\cite{paper1}. 

The task is to compress identical and idependent copies of a density matrix, $\rho^{\otimes n}$. We assume that the density matrix $\rho$ is unknown except for its \emph{spectrum}. Let the memory be~$M$. We say a pair of encoder and decoder $(\E,\D)$ is a $(|M|,\delta)$-code if
\begin{equation}\label{eq:error}
   \frac12 \|\rho^{\otimes n}-\D\circ\E(\rho^{\otimes n})\|_1\le\delta\ .
\end{equation} 

The goal is to determine the optimal sequence of $(|M_n|,\delta_n)$-codes that uses the minimal memory size such that the error vanishes asymptotically at large $n$, $\delta_n=o(1)$. In comparison to Schumacher compression, the error is judged without including the purification. Hence, it is not compressing quantum data but rather compressing the description of the density matrix itself. We call this the \emph{quantum minimum description} of $\rho^{\otimes n}$, and the minimal compression size is the \emph{quantum minimum description length} (QMDL). We would like to compute the asymptotic QMDL as an expansion in $n$ up to terms that vanish at large~$n$. 

Let us consider the following heuristics. Let $\Gamma$ be an $\eps$-net that discretizes the unitary orbit of a given non-degenerate spectrum $\vec p$. Let the set of all possible density matrices $\rho$ with the same spectrum $\vec p$ be $\Omega$. It can be embedded in $\mathbb R^{d(d-1)}$ because $d(d-1)$ is the number of degrees of freedom (dimensions) of the orbit.  The volume of $\Omega$, calculated using the Lebesgue measure, is given by the Vandermonde determinant (up to constant factors)
\[
    \mathrm{Vol}(\Omega) \propto\, \prod_{i<j} (x_i-x_j)^2\ .
\]

On the other hand, this volume can be calculated using a Riemann sum: number of cells $|\Gamma|$ in the $\eps$-net times the unit cell volume
\[
    \mathrm{Vol}(\Omega)\propto |\Gamma| \eps^{d(d-1)}\ .
\]
It follows that 
\[
    \log|\Gamma| \simeq d(d-1)\log\eps^{-1} +2\sum_{i<j} |x_i-x_j| +\mathrm{const.}\ .
\]
This quantity is closely related to free entropy in the operator algebra literature~\cite{voiculescu1993analogues,voiculescu1994analogues,voiculescu2002free} as we show in the companion paper~\cite{paper1}. Free entropy counts the number of possible matrices that share similar moments as $\rho$. In our context, this is the number of bits needed to label different states on the unitary orbit given the spectrum up to a 2-norm precision $\eps$, and it should determine the minimal memory size. As we will show, it turns out that the QMDL of $\rho^{\otimes n}$ is not too different from the formula above. If we set $\eps=n^{-1}$ and divide $\log|\Gamma|$ by a factor of two, we obtain the minimum compression size (QMDL) 
\begin{equation}
 |M_n| = \frac12(d^2-d)\log n +\sum_{i<j}\log|x_i-x_j| + \mathrm{const.}+o(1)\ .
\end{equation}
where the constant actually equals $-\sum_{k=0}^{d-1}\log k!$.

More generally, we obtain the QMDL for any $\rho$ such that all positive eigenvalues being non-degenerate.\footnote{We believe that this non-degeneracy assumption can be lifted. The key is to upgrade Theorem~\ref{thm:main} to deal with degeneracy.} Let the rank of $\rho$ be $r$, the QMDL of $\rho^{\otimes n}$ is given by
\begin{equation}\label{eq:result}
        |M_n| = \frac12 r(2d-r-1)\log n + \sum_{1 \le i < j \le r}\log|x_i-x_j| + (d-r)\sum_{i=1}^r \log x_i - \sum_{k=d-r}^{d-1}\log k! + o(1) \ .
\end{equation}

This compression task has been studied by one of us in~\cite{yang2016optimal,yang2018compression} at the leading order~$O(\log n)$, and we now push it to the second order. Practically, the second order rate is less important, but it unveils itself as the free entropy~\cite{voiculescu1993analogues,voiculescu1994analogues,voiculescu2002free}. In the companion paper~\cite{paper1}, we introduce the free entropy and show its equality to~\eqref{eq:result}.

%cloning map.
Our compression scheme is a generalization of the Yang-Chiribella-Hayashi (YCH) scheme for compressing qubits~\cite{yang2016optimal}. The YCH scheme crucially utilizes Werner's cloning map between two the symmetric subspaces~\cite{werner1998optimal}. It turns out that for the YCH scheme to work for qudits, one needs to generalize Werner's cloning map to two arbitrary GL$(d,\mathbb C)$ irreps, representing the input and output Hilbert spaces respectively. In this context, the notion of ``cloning'' generalizes to mapping some $\mu$-irrep of $\rho$, denoted as $\rho_\mu$, to a clone of another $\nu$-irrep of $\rho$, denoted as $\rho_\nu$. We propose such a generalized cloning map $\C_{\mu\to\nu}$ for any two GL$(d,\mathbb C)$ irreps. We show that this cloning has a high fidelity if the irreps $\mu$ and $\nu$ are close.\footnote{Precisely, we mean that their corresponding Young diagrams are close, as measured by the distance $d(\mu,\nu) := \sum_{i=1}^d |\mu_i-\nu_i|$.} This property is essential for our compression to work.

This paper is structured as follows: We first review the YCH scheme for compressing qubits in Section~\ref{sec:review}. We then generalize Werner's cloning map to arbitrary GL$(d,\mathbb C)$ irreps in Section~\ref{sec:cloner}, and it helps us to generalize the YCH scheme to apply for qudits in Section~\ref{sec:direct}. In Section~\ref{sec:converse}, we show that our compression scheme is optimal. As mentioned above, we used a key fidelity estimate of the generalized cloning map in the achievability proof, and we leave its technical proof to Section~\ref{sec:fidelity}. 


\section{Background}\label{sec:background}
{\color{blue}
\subsection{Related works on the quantum minimum description}\label{sec:related_works}
The quantum extension of the minimum description length \cite{rissanen1978modeling} concerns determining the optimal rate of compressing $n$ copies of a (partially) unknown quantum state up to a vanishing error.
Early theoretical work in 2010 by Plesch and Bu\v{z}ek \cite{plesch2010efficient} established that $n$ copies of an identically prepared pure qubit could be exponentially compressed into an $O(\log n)$ qubit space. This theoretical foundation was soon realized experimentally by Rozema et al.~in 2014 \cite{rozema2014quantum}, who successfully compressed three identically prepared pure qubits into two, proving the physical viability of quantum data compression.
For pure states, the compression scheme is essentially to realize a basis transformation to the symmetric subspace, on which pure states are fully supported. Since the dimension of the symmetric subspace scales only polynomially in $n$, an exponential compression (in the number of qubits) is achieved. For mixed states, the problem is more involved, and its algebraic basis echoes the seminal works of Koashi and Imoto \cite{koashi2001compressibility,koashi2001possible}.

A major theoretical breakthrough was the YCH scheme \cite{yang2016optimal} in 2016, which established the ultimate limits for compressing identically prepared mixed states. Although the authors focused on the first order of the compression rate (i.e., the $\log n$ term of the memory cost), the scheme is optimal to the second order for fixed-spectrum state families, as we show later in the present article.
The scheme was later applied to a task of time tracking \cite{yang2018quantum,yang2018compressionfor}, where compression serves as the key routine of a ``quantum stopwatch" that achieves an information-theoretic advantage over its classical counterparts.
For arbitrary-dimensional diagonal states (classical probability distributions), the optimal compression scheme was proposed by Hayashi and Tan in 2017 \cite{hayashi2017minimum}.
An attempt to generalize the compression to higher-dimensional systems (qudits) was made in \cite{yang2016efficient}, but the scheme there could only achieve an $O(\log n)$ rate with a sub-optimal prefactor. 

In 2018, Yang and colleagues designed a compression scheme for qudits \cite{yang2018compression}, determining that each unknown parameter or degree of freedom of the state generally incurs a memory cost of roughly $\frac{1}{2} \log n$ bits or qubits. Ref.~\cite{yang2018compression} also proved the optimality of the scheme up to the first order, as well as an intriguing result: the memory for storing identically prepared qudits can be a hybrid of qubits and classical bits, but cannot be fully classical.
A recent work by one of us considered states generated by shallow quantum circuits, where not only the number of copies $n$ but also the dimension of the system $d$ is treated as an asymptotic parameter.
On the experimental front, advancements in noisy intermediate-scale quantum (NISQ) devices have allowed researchers to implement these optimal compression algorithms on actual hardware \cite{pivoluska2022implementation,xu2024experimental}. 

We conclude the review with an important remark: The scheme in Ref.~\cite{yang2018compression}, which is the state-of-the-art scheme for qudits, was constructed via a series of asymptotic tools, making the evaluation of its compression rate tractable only at the first order.
Hence, the second-order rate for identically prepared qudit compression remains unknown.
In the remaining part of the present article, we will address this by presenting a compression scheme that rigorously generalizes the YCR scheme from qubits to qudits.}



\subsection{Review of Yang-Chiribella-Hayashi scheme for compressing qubits}\label{sec:review}
The quantum minimal description of qubits, as introduced by Yang, Chiribella, and Hayashi (YCH)~\cite{yang2016optimal}, relies on exploiting the permutation symmetry inherent in identically prepared states. Given $n$ copies of an unknown qubit state $\rho = U(g)\rho_0 U(g)^\dagger$, where $g\in\mathrm{U}(2)$ and $\rho_0$ is diagonal with eigenvalues $p$ and $1-p$ (with $p \ge 1/2$), the goal is to compress this state into a smaller quantum memory and later decompress it with vanishing asymptotic error.

The mathematical engine of the YCH scheme is the Schur transform, denoted as $V_\mathrm{Schur}$, which maps the $n$-qubit Hilbert space to the Schur-Weyl basis:
\begin{equation}
    V_\mathrm{Schur}: (\mathbb{C}^2)^{\otimes n} \to \bigoplus_{J=j_{\min}}^{n/2} \h_J \otimes \map{M}_J
\end{equation}
Here, $\h_J$ denotes the $J$-irreducible representation (irrep) of GL$(2,\mathbb{C})$\footnote{Note that $\h_J$ is the same irrep space for both $\mathrm{GL}(2,\mathbb{C})$ and $\mathrm{U}(2)$.} with dimension $2J+1$, and $\map{M}_J$ denotes the corresponding irrep of the symmetric group $S_n$. For qubits, the allowed irreps are labeled by the integer or half-integer $J$, corresponding to the partition $(n/2+J, n/2-J)$. Here, $j_{\min} = 0$ for even $n$ and $j_{\min} = 1/2$ for odd $n$, which ensures the corresponding partition $(n/2+J, n/2-J)$ has non-negative row lengths with $n/2+J \ge n/2-J$.

Applying the Schur transform to the input state block-diagonalizes it:
\begin{equation}\label{eq:schur}
    V_\mathrm{Schur} \rho^{\otimes n} V_\mathrm{Schur}^\dagger = \bigoplus_{J=j_{\min}}^{n/2} q_J \rho_{g,J} \otimes \tau_J
\end{equation}
where $\rho_{g,J} :=  U_J(g) \rho_J U_J^\dagger(g)$ is the state on the representation space $\h_J$, and $\tau_J$ is the maximally mixed state on the multiplicity space $\map{M}_J$. The dimension of $\map{M}_J$ is given by $(2J+1)\binom{n+1}{n/2+J+1}/(n+1)$. The probability distribution $q_J$ for measuring a specific $J$ sector is defined by the binomial distribution differences:
\[
    q_J :=  \frac{2J+1}{2J}\left[B(n/2+J+1;p,n)-B(n/2-J;p,n)\right]
\]
where $B(k;p,n) :=  \binom{n}{k}p^{k}(1-p)^{n-k}$.\\

With the basis established, the compression scheme proceeds as follows:
\begin{enumerate}
    \item \textbf{Irrep Measurement:} Perform a non-demolition measurement of the quantum number $J$, preserving the quantum information within the associated subspace $\h_J \otimes \map{M}_J$. 
    \item \textbf{Cloning:} Let $J^\star$ be the most typical outcome (the value that maximizes $q_J$), ($J^\star:= (p-1/2)(n+1)$). Discard the multiplicity register $\map{M}_J$ and apply the cloning channel $\C_{J\to J^\star}$ to map the remaining state $\rho_{g,J}$ to a fixed irrep $J^\star$. Store this output state into a quantum memory of dimension $2J^\star + 1$.
\end{enumerate}

To decompress the state, the decoder executes the process in reverse:
\begin{enumerate}
    \item \textbf{Irrep Sampling:} Pick a target value $K$ at random from the distribution $q_K$.
    \item \textbf{Cloning:} Apply the optimal cloning channel $\mathcal{C}_{J^\star\to K}$ to the state in the quantum memory.
    \item \textbf{State Reconstruction:} Append a new multiplicity register initialized in the maximally mixed state $\tau_K$, and perform the inverse Schur transform.
\end{enumerate}

In summary, the YCH scheme works because all the information about $\rho^{\otimes n}$ itself is contained in the GL$(2,\mathbb C)$ irrep. We can discard the maximally mixed state on the symmetric irreps because we do not care about preserving the purification. Furthermore, we want to only keep $\mathcal U_{J^\star}$, as $q_J$ concentrates at $J^\star$. The key is to apply the channel $\C_{J\to J^\star}$ to whatever quantum state $\rho_{g,J}$ measured, so we do not need to spend extra memory to record which $J$ is obtained in the measurement.

A crucial ingredient is the quantum channel $\C_{J\to J^\star}$. It is chosen to be \emph{Werner's cloning map}~\cite{werner1998optimal}. While the fundamental no-cloning theorem prohibits the exact duplication of an unknown quantum state, Werner's cloner provides the optimal approximation for cloning $n$ identical copies of any pure state into $m$ copies ($m > n$) with the maximal fidelity. Werner's cloner is defined as
\begin{equation}\label{eq:werner}
    \mathcal{C}_{n \to m}:\mathrm{Sym}^n(\mathbb{C}^d)\to\mathrm{Sym}^m(\mathbb{C}^d),\quad \rho\mapsto \frac{d_{\mathrm{Sym}^n(\mathbb{C}^d)}}{d_{\mathrm{Sym}^m(\mathbb{C}^d)}}V^\dagger(\rho\otimes \id_{n-m})V
\end{equation}
where $\id_{m-n}$ is the identity on $\mathrm{Sym}^{|n-m|}(\mathbb{C}^d)$, and $V$ is the isometric embedding of $\mathrm{Sym}^m(\mathbb{C}^d)$ into $\mathrm{Sym}^n(\mathbb{C}^d)\otimes \mathrm{Sym}^{|n-m|}(\mathbb{C}^d)$,\footnote{The cloning map is usually written as $\mathcal{C}_{n \to m}(\rho) = \frac{d_{\mathrm{Sym}^n(\mathbb{C}^d)}}{d_{\mathrm{Sym}^m(\mathbb{C}^d)}}P_{\rm{Sym}^m}(\rho\otimes \id_{n-m})P_{\rm{Sym}^m}$, where $P_{\rm{Sym}^m}$ is the projection onto $\mathrm{Sym}^m(\mathbb{C}^d)$ acting on the larger space $\mathrm{Sym}^n(\mathbb{C}^d)\otimes \mathrm{Sym}^{|n-m|}(\mathbb{C}^d)$. We take the intrinsic view by using the isometry $V$ instead, ensuring the channel's codomain is strictly $\mathrm{Sym}^m(\mathbb{C}^d)$.} that is fixed up to a phase by the covariance condition $V U(g)^{\otimes m} = U(g)^{\otimes m}V$ for all $g \in \mathrm{U}(d)$.

The symmetric subspace $\mathrm{Sym}^n(\mathbb{C}^d)$ is spanned by $n$-fold tensor power of pure qudits states, $\mathrm{Sym}^n(\mathbb{C}^d)\equiv\rm{span}\{\ketbra{\psi}{\psi}^{\otimes n}\}_{\ket{\psi}\in\mathbb C^d}$. The cloning map is thus commonly described as cloning $n$ copies of pure qudits to $m$ approximate clones. For pure qubits, the map reads
\begin{equation}
    \C_{n\to m}(\ketbra{\psi}{\psi}^{\otimes n}) =
        \frac{n+1}{m+1} V^\dagger(\ketbra{\psi}{\psi}^{\otimes n} \otimes \id_{m-n)})V,\quad \ket{\psi}\in\mathbb C^2
\end{equation}
where $V$ is the isometric embedding onto $\mathrm{Sym}^m(\mathbb{C}^2)$.

In the context of the Schur transform, the representation space $\h_J$ (which has dimension $2J+1$) is naturally isomorphic as a vector space to the symmetric subspace of $2J$ qubits, $\mathrm{Sym}^{2J}(\mathbb{C}^2)$, carrying the spin-$J$ representation of $\mathrm{SU}(2)$. Because quantum states are invariant under global phases, the adjoint action of $\mathrm{U}(2)$ on this space is identical to that of $\mathrm{SU}(2)$. Consequently, transitioning from an irrep $J$ to a larger irrep $K$ is mathematically equivalent to Werner's optimal cloning from $2J$ to $2K$ qubits. This covariant map (extended to include appropriate partial traces for contractions when $J > K$) is formally defined as:
\begin{equation}\label{eq:qubitcloning}
    \mathcal{C}_{J\to K}(\rho_J) :=  \begin{cases} 
        \frac{2J+1}{2K+1} V^\dagger(\rho_J \otimes \id_{2K-2J})V, & J \le K \\ 
        \mathrm{tr}_{2J-2K}(\rho_J), & J > K 
    \end{cases}\ .
\end{equation}

A defining and indispensable feature of this map is its U$(2)$-covariance. Because isometry commutes with the unitary $U(g)^{\otimes 2K}$, and the identity operator $\id_{2K-2J}$ is trivially invariant under $U(g)^{\otimes 2K-2J}$, the channel is $\mathrm{U}(2)$-covariant. This ensures that the rotational symmetries of the encoded state are perfectly preserved during the mapping to the typical $J^\star$ block:
\begin{equation}
    \mathcal{C}_{J\to K}\left(U_J(g) \rho_J U_J^\dagger(g)\right) = U_K(g) \mathcal{C}_{J\to K}(\rho_J) U_K^\dagger(g)\ .
\end{equation}

YCH demonstrated that this compression scheme achieves an asymptotically vanishing error rate $\delta \sim O(n^{-1/2})$. Because the typical value of $J$ peaks strongly at $J^\star = (p-1/2)(n+1)$, the achievable compression rate at large $n$ scales logarithmically:
\begin{equation}\label{eq:achievable_qubit}
    |M_n|=\log (2J^\star+1) = \log n + \log (2p-1) + o(1)\ .
\end{equation}


\section{Generalized Cloning Map}\label{sec:cloner}

We want to generalize the YCH scheme to qudits. For a qudit in state $\rho$, the Schur-Weyl duality states that
\begin{equation}
    V_\mathrm{Schur}: (\mathbb C^d)^{\otimes n}\to\bigoplus_\lambda\h_\lambda\otimes\map{M}_\lambda
\end{equation}
where $\lambda$ is an \emph{ordered partition} of $n$, ($\lambda\vdash n, \lambda_1\ge\cdots\ge\lambda_d $), which is also known as a \emph{Young diagram}; $\h_\lambda$ denotes the $\lambda$-irrep of GL$(d,\mathbb{C})$; and $\map{M}_\lambda$ denotes the corresponding irrep of the symmetric group $S_n$.

For a density matrix $\rho$ on $\mathbb{C}^d$, we can transform it into the Schur basis:
\begin{equation}\label{eq:schur_d}
    V_\mathrm{Schur}\rho^{\otimes n}V_\mathrm{Schur}^\dagger= \bigoplus_\lambda q_\lambda \rho_\lambda\otimes \tau_\lambda=\sum_\lambda q_\lambda\, \rho_\lambda\otimes \tau_\lambda\otimes\ketbra{\lambda}{\lambda}_\Lambda\ ,\quad
\end{equation}
where $\rho_\lambda$ is the normalized $\mathrm{GL}(d,\mathbb{C})$ irreps of $\rho$, and $\tau_\lambda$ is the maximally mixed state on $\map{M}_\lambda$. We rewrite the direct sum as a usual sum of orthogonal blocks with an attached classical register $\Lambda$ to record the irrep labels. 

Following YCH, we measure the $\Lambda$ register and obtain some outcome $\lambda$. The key is to find a $\mathrm{U}(d)$-covariant generalization of the cloning channel that maps any measured Young diagram $\lambda$ to the most typical one $\lambda^\star:=\mathrm{argmax}_\lambda q_\lambda$.\\


\subsection{Generalized cloner as the minimal covariant quantum channel}\label{sec:definition}
We start by examining covariant quantum channels. We say that a quantum channel $\map{N}:\mathcal B(\h_\mu)\to \mathcal B(\h_\nu)$ is $\mathrm{U}(d)$-covariant (equivariant) with respect to two representations $\mu$ and $\nu$ if 
\begin{equation}\label{eq:covariance}
    \map{N}\circ\map{U}_\mu(g) = \map{U}_\nu(g)\circ\map{N},\quad\forall g\in\mathrm{U}(d)\ ,
\end{equation}
where $\map{U}_\mu(g):= U_\mu(g)(\, \cdot\, )U_\mu^\dagger(g)$ and similarly for $\map{U}_\nu(g)$, and the unitaries $U_\mu(g)$ and $U_\nu(g)$ are representations of $\mathrm{U}(d)$. We are primarily interested in the case where the input space $\h_\mu$ and output space $\h_\nu$ are irreps. (cf. Figure~\ref{fig:covariance_diagram}.)

\begin{figure}
    \centering
    \begin{tikzcd}[row sep=large, column sep=large]
        \mathcal{B}(\mathcal{H}_\mu) \arrow[r, "\mathcal{N}"] \arrow[d, "\mathcal{U}_\mu(g)"'] & \mathcal{B}(\mathcal{H}_\nu) \arrow[d, "\mathcal{U}_\nu(g)"] \\
        \mathcal{B}(\mathcal{H}_\mu) \arrow[r, "\mathcal{N}"'] & \mathcal{B}(\mathcal{H}_\nu)
    \end{tikzcd}
    \caption{Commutative diagram illustrating the $\mathrm{U}(d)$-covariance of a $\mathrm{U}(d)$-covariant map. Applying the group action $\mathcal{U}_\mu(g)(\cdot) = U_\mu(g)(\cdot)U_\mu^\dagger(g)$ prior to the channel yields the identical state as applying the channel followed by the output group action $\mathcal{U}_\nu(g)$.}
    \label{fig:covariance_diagram}
\end{figure}

A useful characterization of channel covariance in terms of the Choi operator $J_{\mathcal{N}}$ reads (cf. for instance, Lemma 11 in~\cite{gschwendtner2021programmability}):
\begin{equation}
    [J_{\mathcal{N}},U_\mu^*\otimes U_\nu]=0\ .
\end{equation}
Here, the relevant representation is the dual (contragredient) representation $\mu^*$ of $\mu$ tensored with $\nu$. We can decompose $\mu^*\otimes\nu$ into irreps of $\mathrm{U}(d)$:
\begin{equation}\label{eq:decomp}
    \mu^*\otimes\nu \simeq \bigoplus_\lambda \lambda\ .
\end{equation}
Schur's lemma implies that $J_\mathcal{N}$ takes the form of a convex combination of projectors $\{\Pi_\lambda\}$ onto these irreps:
\begin{equation}\label{eq:covChoi}
    J_\mathcal{N}\simeq \bigoplus_\lambda c_\lambda \frac{d_\mu}{d_\lambda} \Pi_\lambda\ , \quad \sum_\lambda c_\lambda =1, \quad c_\lambda\ge 0\ ,
\end{equation}
where the factor $d_\mu/d_\lambda$ ensures that the convex combination of Choi operators is properly normalized, satisfying $\Tr\, [(d_\mu/d_\lambda)\Pi_\lambda] = d_\mu$.

Because we are interested in a quantum channel that preserves as much information of $\rho$ as possible, we naturally seek the ``minimal'' Choi operator. That is, we pick the projector onto the \emph{minimal irrep} in the above decomposition, which we denote as $\omega$. Minimality here is defined with respect to the dominance (majorization) partial order among the partitions corresponding to the irreps:
\begin{equation}\label{eq:dominance}
    \lambda\prec\eta\ \iff\  \sum_{i=1}^n \lambda_i\le \sum_{i=1}^n\eta_i,\quad\forall\ 1\le n\le d
\end{equation}
where we label the irreps including multiplicities.

While it is obvious that there is a maximal element $\mu^*+\nu$, there is in fact also a unique (multiplicity=1) minimum $\omega$ that is dominated by every other irrep $\lambda$ in the Littlewood-Richardson decomposition of $\mu^*\otimes\nu$. This is known as the \emph{Parthasarathy-Ranga Rao-Varadarajan (PRV) component}~\cite{parthasarathy1967,kumar1988proof}:
\begin{equation}
    \omega = (\mu^*+w_0(\nu))^+
\end{equation}
where $w_0$ is the longest Weyl group element that reflects the simple roots ($w_0:\alpha_i\mapsto -\alpha_{n-i}$), and the $(\lambda)^+$  denotes the projection onto the fundamental Weyl chamber to obtain a dominant weight if $\lambda$ fails to be dominant, (dominance means $\lambda_1\ge\lambda_2\ge\cdots\ge\lambda_d$.) The irrep $\omega$ is thus defined by its highest weight, $(\mu^*+w_0(\nu))^+$.

We can equivalently express the dual as $\mu^* = - w_0(\mu)$, giving us:
\begin{equation}
    \omega = w_0(\nu-\mu)^+ = (\nu-\mu)^+\ .
\end{equation}
In terms of Young diagrams (highest weight), $\omega$ has the highest weight that is simply the difference $\nu-\mu$ re-ordered as a dominant weight. We will later focus on the case when $\nu$ strictly dominates $\mu$: $$\omega =(\nu-\mu)^+\equiv\nu-\mu\ .$$

We therefore define our \emph{generalized cloning map} (or \emph{generalized cloner}) as follows.\JW{Alternative names: Irrep Adapter, Irrep Converter, PRV channel, ...}

\begin{defn}[Generalized Cloning Map]
Let $\mu$ and $\nu$ be $\mathrm{GL}(d,\mathbb{C})$ irreps, and let $\omega = (\nu-\mu)^+$ be the PRV component in the decomposition of $\mu^*\otimes\nu$. The generalized covariant cloning map $\mathcal{C}_{\mu\to\nu}:\mathcal B(\mathcal{H}_\mu)\to \mathcal B(\mathcal{H}_\nu)$ is defined as
\begin{equation}\label{eq:gen_cloner}
    \mathcal{C}_{\mu\to\nu}: \sigma_\mu\mapsto \Tr_\mu\,\left(J_\omega\cdot\sigma_\mu^\top\otimes \id_\nu\right)\ ,\quad J_\omega:=\frac{d_\mu}{d_\omega} \Pi_{\omega}\ .
\end{equation}
\end{defn}

The map $\mathcal{C}_{\mu\to\nu}$ is completely positive because its Choi operator $J_\omega$ is proportional to a projector. To prove the channel is trace-preserving, we can equivalently show that the partial trace of its Choi operator over the output space equals the identity operator on the input space, i.e., $\Tr_\nu(J_\omega) = \id_{\mu^*}$.

First, consider the partial trace of the projector $\Pi_\omega$. Because $\Pi_\omega$ is the projector onto an invariant subspace of $\mu^* \otimes \nu$, it commutes with the joint group action $U_{\mu^*}(g) \otimes U_\nu(g)$ for all $g \in \mathrm{U}(d)$. Taking the partial trace over $\nu$ preserves covariance on the remaining space, meaning $\Tr_\nu(\Pi_\omega)$ must commute with $U_{\mu^*}(g)$ for all $g$. Since $\mu^*$ is an irreducible representation, Schur's Lemma dictates that any operator commuting with the group action must be proportional to the identity:
\[
    \Tr_\nu(\Pi_\omega) = c\, \id_{\mu^*}\ .
\]
We fix the proportionality constant $c$ by taking the trace of both sides. Recognizing that $\Tr(\Pi_\omega) = d_\omega$ and $\Tr(\id_{\mu^*}) = d_{\mu^*} = d_\mu$, we find:
\[
    d_\omega = c\, d_\mu \implies c = \frac{d_\omega}{d_\mu}\ .
\]
Substituting this back into the expression for our Choi operator $J_\omega$, we get:
\[
    \Tr_\nu(J_\omega) = \frac{d_\mu}{d_\omega} \Tr_\nu(\Pi_\omega) = \frac{d_\mu}{d_\omega} \left( \frac{d_\omega}{d_\mu} \id_{\mu^*} \right) = \id_{\mu^*}\ .
\]
This confirms that the channel is trace-preserving.

Werner's pure state cloning map is a special case of cloning between two symmetric representations (see below). Outside the symmetric pure-state case, however, it is more precisely a covariant irrep-orbit conversion channel rather than literal cloning of $\rho^{\otimes n}$. We view the generalized cloning map as designed to clone a $\mathrm{GL}(d,\mathbb C)$ irrep of any density matrix to another irrep. Nonetheless, it remains an interesting question whether this generalized cloning map can help achieve optimal mixed-state cloning~\cite{scarani2005quantum}. 

\subsection{Properties of the generalized cloner}\label{sec:properties}
We now discuss some properties of the generalized cloning map. They help to justify our proposal as a sound generalization of Werner's cloning map.

\subsubsection{Sanity Checks}
We start with some quick sanity checks. When $\mu=\nu$, $\omega$ is the trivial irrep, and the projector onto the trivial irrep correctly corresponds to the identity channel, which is the best cloning map possible. 

We can also examine the qubit case where the optimal cloner is known~\cite{werner1998optimal}. Here, all irreps are isomorphic to the symmetric subspaces $\mathrm{Sym}^k(\mathbb{C}^2)$ for some $k$. 
\begin{rem}
    Let $\h_\mu=\mathrm{Sym}^n(\mathbb{C}^2)$ and $\h_\nu= \mathrm{Sym}^m(\mathbb{C}^2)$. The generalized cloning map reduces to Werner's cloning map~\eqref{eq:werner}. 
\end{rem}
The irrep decomposition of $\mu^*\otimes\nu$ reads:
\begin{equation}
    \mu^*\otimes\nu \simeq \bigoplus_{k=0}^{\min(m,n)} \mathrm{Sym}^{|n-m+2k|}(\mathbb{C}^2)\ .
\end{equation}
The minimal irrep is $\mathrm{Sym}^{|n-m|}(\mathbb{C}^2)$, corresponding to the highest weight $(\nu-\mu)^+=(|n-m|,0)$. The Kraus rank of this generalized cloning channel is $|n-m|+1$, which uniquely matches the covariant channel used by YCH in~\eqref{eq:qubitcloning}. Note that when $m<n$, this generalized channel naturally reduces to a partial trace.

In fact, when $\omega=\nu-\mu$, we can transition from the Choi representation to a Stinespring representation that makes the connection to Werner's cloning map explicit~\eqref{eq:werner}. To do this, we must rigorously define the isometric embedding $V: \mathcal{H}_\nu \to \mathcal{H}_\mu \otimes \mathcal{H}_\omega$. We fix $V$ up to this phase via the covariance condition, demanding that it intertwines the group actions:\footnote{This reduces to the covariance condition we demand for the isometry in Werner's cloning map. For $\mu=\mathrm{Sym}^n(\h), \nu=\mathrm{Sym}^m(\h)$, we have $\omega=\mathrm{Sym}^{m-n}(\h)$ and $U_\nu(g)=U_\mu(g)\otimes U_\omega(g)=U(g)^{\otimes m}$.}
\begin{equation}
    V U_\nu(g) = (U_\mu(g) \otimes U_\omega(g)) V \quad \text{for all } g \in \mathrm{U}(d) \ .
\end{equation}
Using this isometry, the generalized cloning map can be written as:
\begin{equation}
    \mathcal{C}_{\mu \to \nu}(\sigma_\mu) = \frac{d_\mu}{d_\nu} V^\dagger (\sigma_\mu \otimes \id_\omega) V\ .
\end{equation}
We then also refer to $\omega$ as the \emph{ancilla} system.

We rigorously justify this representation in the following lemma.

\begin{prop}[Stinespring Representation]\label{prop:choi}
    Let $V: \mathcal{H}_\nu \to \mathcal{H}_\mu \otimes \mathcal{H}_\omega$ be the isometric embedding of the irrep $\nu$ onto the tensor product of irrep $\mu$ and $\omega$, and these highest weights satisfy $\nu = \mu + \omega$. The covariant channel $$\mathcal{C}_{\mu \to \nu}(\sigma_\mu) = \frac{d_\mu}{d_\nu} V^\dagger (\sigma_\mu\otimes \id_\omega) V$$ has a Choi matrix $J$ that is proportional to the projector $\Pi_\omega$ onto the PRV component $(\nu - \mu)^+$ in $\nu \otimes \mu^*$. Specifically, $J = \frac{d_\mu}{d_\omega} \Pi_\omega$.
\end{prop}
\begin{proof}See Appendix~\ref{app:proofs}.
\end{proof}


\subsubsection{Minimal or Smallest}
Lemma~\ref{prop:choi} is the compelling evidence that we are on the right track. Another plausible variant of the generalized cloning map is defined via choosing the \emph{smallest} irrep based on its dimension rather than the dominance ordering. It is perhaps a more intuitive choice. However, the smallest irrep does not easily admit a simple algebraic construction, making it technically harder to work with, even if it could also serve our purposes. Interestingly, the smallest irrep is generally different from the minimal irrep (PRV component). 
\begin{rem}[Minimal $\neq$ Smallest]\label{rem:prv_not_dim_min}
The PRV component in \(\mu^*\otimes \nu\) is not, in general, the irreducible summand of smallest dimension. This failure can occur even in the dominant case $\nu=\mu+\omega$. For example, let $\mu=(1,0,0)$ and $\nu=(5,1,0)$ be GL\((3)\) irreps. The PRV component is $\omega:=(\nu-\mu)^+ = \nu-\mu=(4,1,0)$, which has dimension $24$. However, $(5,0,0)$ is also a summand in the decomposition of $\mu^*\otimes \nu$\footnote{This is because $(1,0,0)\otimes (5,0,0)\cong(6,0,0)\oplus (5,1,0)$ from the Pieri rule.} and it is the smallest summand of dimension $21$.
\end{rem}



Furthermore, the generalized cloning map has already found applications in other information processing tasks that rely on covariant channel constructions. For instance, in quantum purity amplification~\cite{keyl2001rate,li2025purity}, the optimal protocol utilizes the generalized cloning map from any irrep $\mu$ to the fundamental irrep, denoted as~$\Box$. The Choi matrix of this map is proportional to the projection onto the minimal irrep $\omega\subset \mu^*\otimes\Box$, rather than the irrep with the smallest dimension. This provides a strong justification for using the minimal (PRV) irrep for generalized cloning maps between arbitrary irreps.


\subsubsection{Reverse cloner}
For a given pair of irreps $(\mu,\nu)$, we are interested in both the forward $\mu\to\nu$ direction under the map $\C_{\mu\to\nu}$ as well as the reverse $\nu\to\mu$ direction under the map $\C_{\nu\to\mu}$.  It turns out that $\C_{\nu\to\mu}$ is the Petz recovery map of $\C_{\mu\to\nu}$ and vice versa. 
\begin{prop}[Reverse Cloner]\label{prop:reverse}
Let $\C_{\mu\to\nu}:\mathcal B(\mathcal H_\mu)\to \mathcal B(\mathcal H_\nu)$
be the generalizer cloning map, and $J_{\mu\to\nu}=\frac{d_\mu}{d_\omega}\Pi_\omega$ be its Choi operator. Let $\sigma_\mu :=  \id_\mu/d_\mu$.
Then the Petz recovery map of \(\C_{\mu\to\nu}\) with respect to \(\sigma_\mu\) is exactly the reverse  cloning map:
\[
\mathcal R_{\sigma_\mu,\C_{\mu\to\nu}}=\C_{\nu\to\mu}\ .
\]
Equivalently,
\[
\C_{\nu\to\mu}(X)=\frac{d_\nu}{d_\mu}\,\C_{\mu\to\nu}^{\dagger}(X),
\]
where \(\C_{\mu\to\nu}^{\dagger}\) is the adjoint map.
\end{prop}
\begin{proof}
See Appendix~\ref{app:proofs}.
\end{proof}
In the pure state case, the reverse cloner corresponds to tracing out the excess systems.

\subsubsection{Cloning Fidelity}
We need a criterion to assess the quality of the generalized cloning map. We judge it by the quality of cloning a $\nu$-irrep of $\rho$ from a $\mu$-irrep of $\rho$. A good cloner should have a high cloning fidelity, $F(\C_{\mu\to\nu}(\rho_\mu), \rho_\nu)$, whenever the highest weights $\mu$ and $\nu$ are similar. We proves this in the following proposition for the restricted case when $\nu-\mu$ is a dominant weight, i.e., $(\nu-\mu)^+=\nu-\mu$.
\begin{thm}(Cloning Fidelity)\label{thm:main}
    Let $\rho$ be a density matrix on $\mathbb{C}^d$ with rank $r \le d$. We assume its non-zero eigenvalues are strictly non-degenerate, satisfying $x_1 > x_2 > \dots > x_r > 0$. Let $\mu, \nu$ be two ordered partitions labeling $\mathrm{GL}(d,\mathbb C)$ irreps such that $(\nu-\mu)$ is a dominant weight. Their sizes are parameterized by a large integer $n$ such that $|\mu|,|\nu|=\Theta(n)$. Let the distance between the irreps be $d(\mu,\nu) := \sum_{i=1}^d |\mu_i-\nu_i|$. Let $\rho_{\mu}$ and $\rho_\nu$ be the normalized $\mathrm{GL}(d,\mathbb C)$ representations of $\rho$. Let $\C_{\mu\to\nu}$ be the generalized cloning map from $\mu$ to $\nu$. For any arbitrarily small $\epsilon>0$, the cloning fidelity is asymptotically bounded by:
    \begin{equation}
        F(\C_{\mu\to\nu}(\rho_\mu), \rho_\nu)\ge 1-O\left(\frac{d(\mu,\nu)}{n^{1-\epsilon}}\right)\ ,\quad F(\C_{\nu\to\mu}(\rho_\nu), \rho_\mu)\ge 1-O\left(\frac{d(\mu,\nu)}{n^{1-\epsilon}}\right) \ .
    \end{equation}
\end{thm}
We use the unsquared fidelity $F(\rho,\sigma):= \|\sqrt{\rho}\sqrt{\sigma}\|_1$ throughout the paper. This fidelity bound is crucial for the achievability analysis in the following section and will be proven in Section~\ref{sec:fidelity}. 

In fact, the following lemma shows that $\C_{\mu\to\nu}(\rho_\mu)$ commutes with $\rho_\nu$, so their difference is only in the eigenvalues.
\begin{prop}[Commutativity]\label{prop:commutativity}
    Let $\rho \in \mathcal{D}(\mathbb{C}^d)$ be a density matrix. Let $\mu$ and $\nu$ be irreps of $\mathrm{GL}(d,\mathbb C)$, $ \rho_\mu$ and $ \rho_\nu$ be the normalized $\mathrm{GL}(d,\mathbb C)$ irreps of $\rho$.
    Let $\C_{\mu \to \nu}: \mathcal{B}(\h_\mu) \to \mathcal{B}(\h_\nu)$ be a quantum channel that is $\mathrm{U}(d)$-covariant, i.e., for all $g \in \mathrm{U}(d)$:
    \begin{equation}
        \C_{\mu \to \nu}( U_\mu(g) X U_\mu(g)^{-1} ) = U_\nu(g) \C_{\mu \to \nu}(X) U_\nu(g)^{-1}
    \end{equation}
    where $U_\mu(g) := \pi_\mu(g)$ for $g \in \mathrm{U}(d)$.
    Then, we have
    \begin{equation}
        \left[ \C_{\mu \to \nu}(\rho_\mu), \rho_\nu \right] = 0 \ .
    \end{equation}
\end{prop}
\begin{proof}See Appendix~\ref{app:proofs}.\end{proof}

It would be nice to show that among all covariant channels $\mathcal{N}$ mapping from $\mu$ to $\nu$, the generalized cloning map maximizes the cloning fidelity $F(\mathcal{N}(\rho_\mu), \rho_\nu)$. This would help to justify our proposal. We leave this as a conjecture for future investigations. 

\begin{comment}
\subsubsection{Optimality}
\begin{prop}[Optimality of the Generalized Cloning Map]
\label{thm:optimality}
Let $\mu$ and $\nu$ be irreps of GL$(d,\mathbb C)$, and $\rho_\mu$ and $\rho_\mu$ be corresponding normalized irreps of $\rho$. Among all covariant quantum channels $\mathcal{N}: \mathcal{B}(\mathcal{H}_\mu) \to \mathcal{B}(\mathcal{H}_\nu)$, the generalized cloning map $\mathcal{C}_{\mu \to \nu}$ defined by its Choi operator $J_\omega = \frac{d_\mu}{d_\omega} \Pi_\omega$, where $\omega := (\nu - \mu)^+$, maximizes the overlap $\Tr\,(\mathcal{N}(\rho_\mu)\rho_\nu)$.
\end{prop}
\begin{proof}See Appendix~\ref{app:proofs}.
\end{proof}

We instead provide a different perspective on why the generalized cloner is a natural choice for our task. The PRV sector is the least excited, as measured by the Casimir.
\begin{prop}[Variational Characterization]\label{prop:casimir_prv}
Let $\mathcal N:\mathcal B(\mathcal H_\mu)\to \mathcal B(\mathcal H_\nu)$
be a \(\mathrm U(d)\)-covariant quantum channel, and let \(J_\mathcal N\) be its Choi operator. Let
$\omega:=(\nu-\mu)^+$ be the PRV component in \(\mu^*\otimes \nu\), and let \(C_2\) denote the quadratic Casimir operator of the representation \(U_\mu^*\otimes U_\nu\) acting on \(\mathcal H_\mu^*\otimes \mathcal H_\nu\). Define the average Casimir of the channel as
\begin{equation}
K(\mathcal N):= \frac{1}{d_\mu}\Tr(C_2 J_\mathcal N),
\qquad
d_\mu:= \dim \mathcal H_\mu.
\end{equation}
Then
\begin{equation}
    K(\mathcal N)\ge C_2(\omega),
\end{equation}
with equality if and only if $N=\C_{\mu\to\nu}$.
Equivalently, the PRV channel is the unique \(\mathrm U(d)\)-covariant channel of minimal average Casimir.
\end{prop}

\begin{proof}
    See Appendix~\ref{app:proofs}.
\end{proof}

The functional \(K(\mathcal N)\) should be viewed as a channel-level measure of how high in the irreducible tower of \(\mu^*\otimes \nu\) the Choi weight of \(\mathcal N\) sits. The PRV channel is singled out because it is the unique covariant channel whose Choi operator lies entirely in the lowest possible irreducible sector. In this sense, the PRV channel is the covariant channel of least representation-theoretic excitation. In Lemma~\ref{lem:perturbation}, we show that this least-excitation property is conceptually what yields the highest-weight alignment and shallow-weight stability needed for high fidelity cloning as in Theorem~\ref{thm:main}.
\end{comment}

\section{Achievability}\label{sec:direct}

We now show that the YCH scheme equipped with the generalized cloning map also works for qudits. 
The compression scheme for qudits proceeds as follows. 

Encoding:
\begin{enumerate}
    \item \textbf{Irrep Measurement:} Perform a non-demolition measurement of the representation label $\lambda$ (the Young diagram), preserving the quantum information within the associated Schur-Weyl subspace $\mathcal{H}_\lambda \otimes \mathcal{M}_\lambda$. 
    \item \textbf{Cloning:} Let $\Lambda^\star$ be the universal target representation, constructed from the expected spectrum with sufficient padding to ensure strict dominance over all typical states, $\Lambda^\star_i = \lceil n x_i + (r-i)\xi_n \rceil$ where $\xi_n:=\sqrt{2n}\log n+1$. Discard the multiplicity register $\mathcal{M}_\lambda$ and apply the generalized cloner $\mathcal{C}_{\lambda\to \Lambda^\star}$ to map the remaining state $\rho_\lambda$ to a fixed irrep $\Lambda^\star$. Store this output state into a quantum memory of dimension $d_{\Lambda^\star}$.
\end{enumerate}

Decoding:
\begin{enumerate}
    \item \textbf{Irrep Sampling:} Pick a target representation $\lambda'$ at random from the empirical Schur-Weyl distribution $q_{\lambda',n}$.
    \item \textbf{Cloning:} Apply the optimal cloning channel $\mathcal{C}_{\Lambda^\star\to \lambda'}$ to the state in the quantum memory.
    \item \textbf{State Reconstruction:} Append a new multiplicity register initialized in the maximally mixed state $\tau_{\lambda'}$, and perform the inverse Schur transform.
\end{enumerate}

The following theorem shows the soundness of the above scheme.
\begin{thm}[Achievability]\label{thm:achievability}
Let $\rho$ be a density matrix with rank $r\le d$ and strictly decreasing positive eigenvalues $x_1>\cdots>x_r>0$. Let $\rho_g:= U(g)\rho U(g)^\dagger$ for $g\in \mathrm U(d)$.
There exists a sequence of compression codes $(\mathcal E_n,\mathcal D_n)$ with quantum memory $M_n$ of size
\[
    |M_n| = \frac{r(2d-r-1)}{2}\log n + \sum_{1 \le i < j \le r}\log|x_i-x_j| + (d-r)\sum_{i=1}^r \log x_i - \sum_{k=d-r}^{d-1}\log k! + o(1)
\]
such that
\[
    \delta_n:= 
    \sup_{g\in \mathrm U(d)}
    \frac12\left\|
        \mathcal D_n\circ \mathcal E_n(\rho_g^{\otimes n}) - \rho_g^{\otimes n}
    \right\|_1
    \xrightarrow[n\to\infty]{}0 .
\]
and $\delta_n=O(n^{-\frac14+\epsilon})$ for any arbitrarily small $\epsilon>0$.
\end{thm}

More explicitly, we show that there exists a compression protocol with the above memory cost and an error 
\begin{equation}
    \delta_n\le O(n^{-\frac14+\epsilon}) + 2(n+1)^{\frac{r(r+1)}{2}}e^{-\Delta_n},
\end{equation}
where $\Delta_n>0$ is a slowly growing function of $n$, which we pick to be $\Delta_n=(\log n)^2$, such that the second term is superpolynomially small.

For pure states, we have
\begin{equation}\label{eq:achievable_pure}
    |M_n|=(d-1)\log n -\log (d-1)!+ o(1)\ .
\end{equation}

Before presenting the proof, we isolate the asymptotic evaluation of the representation dimension into a separate lemma.

\begin{lem}[Asymptotic Weyl Dimension]\label{lem:weyl_asymptotic}
    Let $x$ be a probability distribution over $d$ outcomes with exactly $r \le d$ strictly positive, non-degenerate components $x_1 > \dots > x_r > 0$. Let $\lambda$ be an ordered partition of size $n$ such that $\lambda_i = n x_i + O(\sqrt{n}\Delta_n)$ for $1 \le i \le r$, and $\lambda_i = 0$ for $i > r$, where $\Delta_n = o(\sqrt{n})$. The dimension of the $\mathrm{GL}(d, \mathbb{C})$ irrep $\mathcal{H}_\lambda$ expands in~$n$ as:
        $$\log \dim \mathcal{H}_\lambda = \frac{r(2d-r-1)}{2}\log n + \sum_{1 \le i < j \le r}\log(x_i-x_j) + (d-r)\sum_{i=1}^r \log x_i - \sum_{k=d-r}^{d-1}\log k! + O\left(\frac{\Delta_n}{\sqrt{n}}\right) \ .$$
\end{lem}
\begin{proof}
The exact dimension is given by the Weyl dimension formula:
\begin{equation}\label{eq:weyl_dim}
    \dim \h_\lambda = \prod_{1 \le i < j \le d} \frac{\lambda_i - \lambda_j + j - i}{j - i} \ .
\end{equation}
We split the product in the numerator into three distinct blocks based on the rank $r$:
\begin{itemize}
    \item \textbf{Non-zero vs. Non-zero rows ($1 \le i < j \le r$):} 
    Substituting $\lambda_k = n x_k + O(\sqrt{n}\Delta_n)$, the terms scale as $\lambda_i - \lambda_j = n(x_i - x_j) \left( 1 + O\left(\frac{\Delta_n}{\sqrt{n}}\right) \right)$. Factoring out the leading terms yields:
    \[
        \prod_{1 \le i < j \le r} n(x_i - x_j) \left( 1 + O\left(\frac{\Delta_n}{\sqrt{n}}\right) \right) = n^{\frac{r(r-1)}{2}} \left( \prod_{1 \le i < j \le r} (x_i - x_j) \right) \left( 1 + O\left(\frac{\Delta_n}{\sqrt{n}}\right) \right) \ .
    \]
    
    \item \textbf{Non-zero vs. Empty rows ($1 \le i \le r < j \le d$):}
    Since $\lambda_j = 0$, the terms scale as $\lambda_i + j - i = n x_i \left( 1 + O\left(\frac{\Delta_n}{\sqrt{n}}\right) \right)$. There are exactly $r(d-r)$ such terms, giving:
    \[
        \prod_{i=1}^r \prod_{j=r+1}^d n x_i \left( 1 + O\left(\frac{\Delta_n}{\sqrt{n}}\right) \right) = n^{r(d-r)} \left( \prod_{i=1}^r x_i \right)^{d-r} \left( 1 + O\left(\frac{\Delta_n}{\sqrt{n}}\right) \right) \ .
    \]
    
    \item \textbf{Empty vs. Empty rows ($r < i < j \le d$):}
    Both row lengths are strictly zero, leaving exactly $j - i$. This block evaluates strictly to a product of factorials: $\prod_{r < i < j \le d} (j - i) = \prod_{k=1}^{d-r-1} k!$.
\end{itemize}

The denominator of the original formula evaluates to $\prod_{k=1}^{d-1} k!$. Dividing the exact factorial block from the numerator (Block 3) by the denominator leaves:
\[
    \frac{\prod_{k=1}^{d-r-1} k!}{\prod_{k=1}^{d-1} k!} = \frac{1}{\prod_{k=d-r}^{d-1} k!} \ .
\]

Multiplying the polynomial factors of $n$ from Blocks 1 and 2 gives $n^{r(r-1)/2} n^{r(d-r)} = n^{r(2d-r-1)/2}$. Taking the logarithm of the combined expression yields the final formula, with the multiplicative error becoming an additive shift: $\log(1 + O(\Delta_n/\sqrt{n})) = O(\Delta_n/\sqrt{n})$.
\end{proof}


\begin{proof}[Proof of Theorem~\ref{thm:achievability}]
We shall omit the subscript $g$ in $\rho_g$ in this proof for notational clarity. By Schur-Weyl duality, since $\rho$ has rank $r$, the state $\rho^{\otimes n}$ is strictly supported on irreps with at most $r$ rows.  Thus, it is unitarily equivalent to a block diagonal state
\begin{equation}
    \rho^{\otimes n}\simeq\bigoplus_{\lambda : \ell(\lambda) \le r} q_{\lambda,n}\left(\rho_\lambda\otimes\tau_\lambda\right).
\end{equation}
By introducing $P_\lambda$ as the projection onto the irreducible block associated with the Young diagram $\lambda$, we have
\begin{equation}
    P_\lambda\rho^{\otimes n}P_\lambda = q_{\lambda,n}\left(\rho_\lambda\otimes\tau_\lambda\right).
\end{equation}
The encoder is defined as
\begin{equation}
    \map{E}(\cdot) = \sum_{\lambda : \ell(\lambda) \le r}\map{C}_{\lambda\to \Lambda^\star}\circ\Tr_{\map{M}_{\lambda}}(P_\lambda(\cdot)P_\lambda),
\end{equation}
where $\map{C}_{\lambda\to\Lambda^\star}$ is the generalized covariant map from Theorem~\ref{thm:main}. We define our target representation $\Lambda^\star$ to be an ideal central representation that strictly dominates all typical representations. To ensure the weight difference $\Lambda^\star - \lambda$ is dominant for all typical $\lambda$, we pad the components with the maximum typical fluctuation:
\begin{equation}
    \Lambda^\star_i :=   \lceil n x_i + (r-i)\xi_n \rceil \ , \quad \text{where } \xi_n = 2\sqrt{n\Delta_n/2} + 1 \ .
\end{equation}
Note that $\Lambda^\star$ is not a partition of $n$ but of a slightly larger integer. This specific choice of $\xi_n$ is rigorously necessitated to guarantee that the weight difference $\Lambda^\star - \lambda$ is a dominant weight for every state in the typical set. For the difference to form a valid Young diagram, its row lengths must be non-increasing, which requires the gap condition $(\Lambda^\star_i - \lambda_i) - (\Lambda^\star_{i+1} - \lambda_{i+1}) \ge 0$, or equivalently:
\begin{equation}\label{eq:dominance_gap}
    \Lambda^\star_i - \Lambda^\star_{i+1} \ge \lambda_i - \lambda_{i+1} \ .
\end{equation}
Let $\epsilon_n = \sqrt{n\Delta_n/2}$ denote the maximum typical fluctuation bound from Eq.~\eqref{eq:typical_set}. The right-hand side is maximized when adjacent rows fluctuate in opposite extreme directions, bounded by:
\begin{equation}
    \lambda_i - \lambda_{i+1} \le n(x_i - x_{i+1}) + 2\epsilon_n \ .
\end{equation}

Conversely, using the standard inequality for ceiling functions $\lceil a \rceil - \lceil b \rceil \ge a - b - 1$, the minimum guaranteed gap of our constructed target representation is:
\begin{equation}
    \Lambda^\star_i - \Lambda^\star_{i+1} \ge n(x_i - x_{i+1}) + \xi_n - 1 \ .
\end{equation}
Substituting these worst-case extrema into Eq.~\eqref{eq:dominance_gap}, the macroscopic $n(x_i - x_{i+1})$ terms exactly cancel, leaving the strict requirement:
\begin{equation}
    \xi_n - 1 \ge 2\epsilon_n \ .
\end{equation}
Setting $\xi_n = 2\epsilon_n + 1$ provides the tightest possible sub-extensive buffer that unconditionally absorbs both the simultaneous worst-case statistical fluctuations and the discrete rounding errors. This geometrically ensures that $\Lambda^\star$ serves as a universal dominant umbrella for the entire typical set.


The corresponding decoder maps back to the original blocks:
\begin{equation}
    \map{D}_n(\cdot) = \bigoplus_{\lambda : \ell(\lambda) \le r} q_{\lambda,n}\left(\map{C}_{\Lambda^\star\to \lambda}(\cdot)\otimes\tau_\lambda\right).
\end{equation}
The trace distance error of the compression protocol evaluates to:
\begin{equation}
\begin{aligned}
    \delta_n &= \frac12\left\|\map{D}_n\circ\map{E}_n(\rho^{\otimes n})-\rho^{\otimes n}\right\|_1\\
    &= \frac12\left\|\bigoplus_{\lambda}q_{\lambda,n}\left(\map{C}_{\Lambda^\star\to \lambda}\circ\map{C}_{\lambda\to \Lambda^\star}(\rho_\lambda)-\rho_\lambda\right)\otimes\tau_\lambda\right\|_1\\
    &= \sum_{\lambda}q_{\lambda,n} \frac12 \left\|\map{C}_{\Lambda^\star\to \lambda}\circ\map{C}_{\lambda\to \Lambda^\star}(\rho_\lambda)-\rho_\lambda\right\|_1.
\end{aligned}
\end{equation}
Because the empirical distribution $\lambda/n$ is restricted to $r$ non-zero components, we define the typical subset of Young diagrams as:
\begin{equation}\label{eq:typical_set}
\mathcal{T}_{p,n}:= \left\{\lambda~:~\lambda_i = 0 \text{ for } i>r, \ \max_{1 \le i \le r}|\lambda_i-x_i n|\le \sqrt{n\Delta_n/2} \right\}.
\end{equation}
The error splits into typical and atypical contributions:
\begin{equation}\label{eq:error_twoterms}
\delta_n\le \max_{\lambda\in\mathcal{T}_{p,n}} \frac12\left\|\map{C}_{\Lambda^\star\to \lambda}\circ\map{C}_{\lambda\to \Lambda^\star}(\rho_\lambda)-\rho_\lambda\right\|_1 + \delta_{\rm tail}(q_{\lambda,n}) \ , \quad \delta_{\rm tail}(q_{\lambda,n}):= \sum_{\lambda\not\in\mathcal{T}_{p,n}}q_{\lambda,n}, 
\end{equation}
having used the maximal trace distance bound $\frac12\|\rho-\sigma\|_1\le 1$ for the atypical term.

The tail mass $\delta_{\rm tail}$ is controlled by Sanov's theorem via the following lemma:
\begin{lem}[Sanov's theorem~\cite{yang2016efficient,christandl2006spectra,o2021quantum,hayashi2017group}]\label{lem:tail_prob}
    Let $\lambda/n$ be the empirical probability distribution associated with the Young diagram $\lambda$.
    The probability $\mathsf{Pr}[\lambda:d(\lambda/n,p)>x]$ is upper bounded by $(n+1)^{r(r+1)/2} e^{-2nx^2}$, where $d(\lambda/n,p):= (1/2)\sum_{i=1}^d|\lambda_i/n-x_i|$ is the total variation distance.
\end{lem}
For any $\lambda\not\in\mathcal{T}_{p,n}$, there exists a $\lambda_i$ such that $|\lambda_i/n-x_i|>\sqrt{\Delta_n/(2n)}$. Thus, the total variation distance is strictly bounded below by $\sqrt{\Delta_n/(2n)}$, and Lemma \ref{lem:tail_prob} yields:
\begin{equation}\label{eq:tail_prob_bound}
    \delta_{\rm tail}(q_{\lambda,n})\le (n+1)^{r(r+1)/2} e^{-\Delta_n}.
\end{equation} 

For typical states $\lambda\in\mathcal{T}_{p,n}$, we bound the trace distance using the triangle inequality and the contractivity of trace distance under CPTP maps:
\begin{equation}\label{eq:error_bound1}
\begin{aligned}
    \frac12\left\|\map{C}_{\Lambda^\star\to \lambda}\left(\map{C}_{\lambda\to \Lambda^\star}(\rho_\lambda)\right)-\rho_\lambda\right\|_1 
    &\le \frac12\left\|\map{C}_{\Lambda^\star\to \lambda}\left(\map{C}_{\lambda\to \Lambda^\star}(\rho_\lambda)\right) - \map{C}_{\Lambda^\star\to \lambda}(\rho_{\Lambda^\star})\right\|_1 + \frac12\left\|\map{C}_{\Lambda^\star\to \lambda}(\rho_{\Lambda^\star}) - \rho_\lambda\right\|_1 \\
    &\le \frac12\left\|\map{C}_{\lambda\to \Lambda^\star}(\rho_\lambda) - \rho_{\Lambda^\star}\right\|_1 + \frac12\left\|\map{C}_{\Lambda^\star\to \lambda}(\rho_{\Lambda^\star}) - \rho_\lambda\right\|_1 \ .
\end{aligned}
\end{equation}

By Theorem~\ref{thm:main}, the fidelity of the generalized cloning map is controlled by the distance between the irreps. For any typical state $\lambda \in \mathcal{T}_{p,n}$, the distance to the target umbrella state $\Lambda^\star$ evaluates to:
\begin{equation}
    d(\Lambda^\star, \lambda) = \sum_{i=1}^r |\Lambda^\star_i - \lambda_i| \le \sum_{i=1}^r \left( |\lceil n x_i + (r-i)\xi_n \rceil - n x_i| + |n x_i - \lambda_i| \right) \ .
\end{equation}
Using the bounds from the typical set definition~\eqref{eq:typical_set} and noting that $\xi_n = O(\sqrt{n\Delta_n})$, this distance is bounded by $O(\sqrt{n\Delta_n})$. 
Choosing $\Delta_n = (\log n)^2$, we guarantee that the distance scales as $d(\Lambda^\star, \lambda) = O(n^{\frac{1}{2}+\epsilon'})$ for any $\epsilon' > 0$. Substituting this distance into Theorem~\ref{thm:main}, the fidelity for both covariant maps is bounded by $1 - O(n^{-\frac{1}{2} + (\epsilon'+\epsilon)})$.

Applying the Fuchs-van de Graaf inequalities, the trace distance error for both maps is universally bounded by:
\begin{equation}
    \frac{1}{2}\|\rho - \sigma\|_1 \le \sqrt{2(1-F)} = \sqrt{ O\left( n^{-\frac{1}{2} + (\epsilon'+\epsilon)} \right) } = O\left( n^{-\frac{1}{4} + \epsilon''} \right) \ ,
\end{equation}
where $\epsilon'' = (\epsilon'+\epsilon)/2$ remains arbitrarily small. Substituting these back into Eq.~(\ref{eq:error_bound1}), the error for typical states is exactly $O(n^{-\frac14+\epsilon''})$. Substituting into Eq.~(\ref{eq:error_twoterms}), the total trace distance error evaluates to:
\begin{equation}
    \delta \le O(n^{-\frac14+\epsilon''}) + (n+1)^{\frac{r(r+1)}2} e^{-\Delta_n} \ .
\end{equation}

Meanwhile, since the input state is mapped strictly into the single $\mathrm{GL}(d,\mathbb C)$ irreducible subspace $\mathcal{H}_{\Lambda^\star}$, the memory cost is exactly $|M_n| = \log d_{\Lambda^\star}$. Since the components of $\Lambda^\star$ explicitly satisfy $\Lambda^\star_i = n x_i + O(\sqrt{n\Delta_n})$, we apply Lemma~\ref{lem:weyl_asymptotic} with sub-extensive fluctuations of $O(\sqrt{n\Delta_n})$. The $O(\sqrt{n\Delta_n})$ shift vanishes entirely into the $o(1)$ tail, yielding exactly:
\[
    |M_n| = \frac{r(2d-r-1)}{2}\log n + \sum_{1 \le i < j \le r}\log(x_i-x_j) + (d-r)\sum_{i=1}^r \log x_i - \sum_{k=d-r}^{d-1}\log k! + o(1) \ .
\]
\end{proof}

\section{Converse}\label{sec:converse}

We now prove the optimality of the compression rate.
Because the source state is unknown except for its spectrum, the converse must be stated uniformly over the full unitary orbit
\[
\rho_g :=  U(g)\rho U(g)^\dagger, \qquad g\in \mathrm U(d).
\]

We first prove an intermediate result.
\begin{prop}[Irreducible orbit sector compression]\label{prop:orbit_sector_converse}
Let $\{\lambda^{(n)}\}_{n\ge 1}$ be any sequence of partitions with length $\ell(\lambda^{(n)})\le d$, and define the orbit ensemble equipped with the Haar measure $\dd g$:
\begin{equation}
    \mathfrak E_{\lambda^{(n)}} :=  \{\dd g,\rho_{g,\lambda^{(n)}}\}_{g\in \mathrm U(d)},
    \qquad
    \rho_{g,\lambda}:= U_\lambda(g)\rho_\lambda U_\lambda(g)^\dagger .
\end{equation}
Let $M_n$ be a quantum memory, and suppose there exists a sequence of sector codes $(\mathcal E_n,\mathcal D_n)$ for the ensemble $\mathfrak E_{\lambda^{(n)}}$ with Haar-average error
\begin{equation}
    \int_{\mathrm U(d)} \dd g\;
    \frac12\left\|
        \mathcal D_n\circ \mathcal E_n(\rho_{g,\lambda^{(n)}})
        - \rho_{g,\lambda^{(n)}}
    \right\|_1
    \xrightarrow[n\to\infty]{}0 .
\end{equation}
Then
\begin{equation}
    |M_n| \ge \log \dim \h_{\lambda^{(n)}} + o(1).
\end{equation}
\end{prop}

\begin{rem}
    In the broader context of our full protocol, a \emph{visible} code means the encoder is told the classical label $\lambda$ (the Schur sector) and can tailor its operation accordingly. However, within any specific sector $\lambda$, the encoder does not know which member of the continuous source family $g \in \mathrm{U}(d)$ was prepared. Therefore, the sector code $(\mathcal E_n,\mathcal D_n)$ analyzed in this proposition is $\lambda$-\emph{visible}, but strictly $g$-\emph{blind}.
\end{rem}

\begin{proof}
Fix $n$ and write $\lambda=\lambda^{(n)}$ for brevity.
Let
\begin{equation}
    \mathfrak A_\lambda
    := 
    C^*\!\left(\,\rho_{g,\lambda}:g\in \mathrm U(d)\,\right)
    \subseteq \mathcal B(\h_\lambda)
\end{equation}
be the unital\footnote{An algebra is unital if contains a multiplicative identity.} $C^*$-algebra generated by the orbit ensemble.

We first show that $\mathfrak A_\lambda = \mathcal B(\h_\lambda)$.
Let $|v_\lambda\rangle\in \h_\lambda$ be a highest-weight vector.
Given the spectrum $x$ of $\rho$, the eigenvalue of $\rho_\lambda$ corresponding to the highest weight is $x^\lambda/s_\lambda(x)$,
and it is strictly larger than every other eigenvalue of $\rho_\lambda$.
Moreover, the highest-weight space is one-dimensional.
Hence the top eigenspace of $\rho_\lambda$ is rank one, and its spectral projector is $\ketbra{v_\lambda}{v_\lambda}$. It can be expressed as a polynomial in $\rho_\lambda$ via the Lagrange interpolating polynomial:
\begin{equation}
\ketbra{v_\lambda}{v_\lambda} = \prod_{j \neq 1} \frac{\rho_\lambda - x_j \id}{x_1 - x_j}
\end{equation}
where $x_i$ denotes the distinct eigenvalues of $\rho_\lambda$ in decreasing order.

Since this expression involves only scalar multiplication, identity operators, and powers of $\rho_\lambda$, it follows that 
\begin{equation}
U_\lambda(g) \ketbra{v_\lambda}{v_\lambda} U_\lambda(g)^\dagger \in \mathfrak{A}_\lambda.
\end{equation}

Now let $X\in \mathfrak A_\lambda'$ belong to the commutant of $\mathfrak A_\lambda$.
Then $X$ commutes with every rank-one projector in the orbit above.
Writing
\begin{equation}
|u_g\rangle:= U_\lambda(g)|v_\lambda\rangle,
\end{equation}
the relation
\begin{equation}
\left[X,\ketbra{u_g}{u_g}\right]=0
\end{equation}
implies that $X|u_g\rangle$ lies in the one-dimensional range of $\ketbra{u_g}{u_g}$.
Hence for each $g$ there exists a scalar $a(g)$ such that
\begin{equation}
    X|u_g\rangle = a(g)\,|u_g\rangle .
\end{equation}
The function $g\mapsto a(g)$ is continuous and takes values in the finite set $\mathrm{spec}(X)$.
Since $\mathrm U(d)$ is connected, $a(g)$ must be constant $a(g):= a$. 
Thus every orbit vector $|u_g\rangle$ is an eigenvector of $X$ with the same eigenvalue $a$.

Finally, the linear span of the orbit
\begin{equation}
\{U_\lambda(g)|v_\lambda\rangle : g\in \mathrm U(d)\}
\end{equation}
is a nonzero $U_\lambda$-invariant subspace of the irrep $\h_\lambda$, and therefore equals all of $\h_\lambda$.
So $X=a\id_{\h_\lambda}$, proving
\begin{equation}
    \mathfrak A_\lambda'=\mathbb C\, \id_{\h_\lambda}.
\end{equation}
By the finite-dimensional bicommutant theorem,
\begin{equation}
    \mathfrak A_\lambda = \mathcal B(\h_\lambda).
\end{equation}

We now invoke the Koashi-Imoto structure theorem for \emph{blind} compression of mixed-state ensembles~\cite{koashi2001compressibility,koashi2001possible}:
the asymptotically necessary quantum memory is the logarithm of the dimension of the nonredundant quantum factor determined by the $C^*$-algebra generated by the ensemble.
Because the generated algebra here is the full matrix algebra $\mathcal B(\h_\lambda)$, the nonredundant quantum factor is the whole irrep space $\h_\lambda$.
Therefore, any $g$-blind code for this sector with vanishing Haar-average error must satisfy
\begin{equation}
    |M_n| \ge \log \dim \h_\lambda + o(1).
\end{equation}
This proves the claim.
\end{proof}

\begin{thm}[Converse]\label{thm:converse}
Let $\rho$ be a density matrix with rank $r\le d$ and spectrum $x$. Let $\rho_g:= U(g)\rho U(g)^\dagger$ for $g\in \mathrm U(d)$.
Suppose $(\mathcal E_n,\mathcal D_n)$ is a sequence of compression codes with quantum memory $M_n$ such that
\begin{equation}
    \delta_n:= 
    \sup_{g\in \mathrm U(d)}
    \frac12\left\|
        \rho_g^{\otimes n}
        -
        \mathcal D_n\circ \mathcal E_n(\rho_g^{\otimes n})
    \right\|_1
    \xrightarrow[n\to\infty]{}0 .
\end{equation}
Then, the required quantum memory size is bounded below by the minimum dimension among the typical Schur sectors:
\begin{equation}
    |M_n| \ge \min_{\lambda \in \mathcal T_{p,n}} \log \dim \h_{\lambda} + o(1),
\end{equation}
where $\mathcal T_{p,n}$ is the typical set of partitions defined by
\begin{equation}
    \mathcal T_{p,n}:= \left\{ \lambda : \ell(\lambda)\le r, \max_{1\le i\le r} |\lambda_i-nx_i| \le \sqrt{n}\log n \right\}.
\end{equation}

Furthermore, if the non-zero eigenvalues are strictly decreasing ($x_1>\cdots>x_r>0$), this evaluates asymptotically to:
\begin{equation}
    |M_n|
    \ge
    \frac{r(2d-r-1)}{2}\log n
    + \sum_{1\le i<j\le r}\log (x_i-x_j)
    + (d-r)\sum_{i=1}^r \log x_i
    - \sum_{k=d-r}^{d-1}\log k!
    + o(1).
\end{equation}
\end{thm}

\begin{proof}
We first outline the core strategy of the proof. Instead of analyzing the compression of the entire ensemble directly, we decompose the input state into its Schur sectors. If a sequence of compression codes achieves vanishing worst-case error overall, probabilistic arguments guarantee that it must also perform well on at least one sequence of typical irreps $\rho_{\lambda^{(n)}}$ where $\lambda^{(n)} \in \mathcal T_{p,n}$. Consequently, the required quantum memory must be bounded below by the dimension of this typical irrep space. When the spectrum is non-degenerate, we apply the Weyl dimension formula to obtain the precise asymptotic expansion.

Applying the Schur transform, the input state takes the form
\begin{equation}\label{eq:converse_schur_form}
    \widehat\rho_{g,n}
    := 
    V_{\mathrm{Schur}}\, \rho_g^{\otimes n}\, V_{\mathrm{Schur}}^\dagger
    =
    \sum_{\lambda:\,\ell(\lambda)\le r}
    q_{\lambda,n}\,
    \rho_{g,\lambda}\otimes \tau_\lambda \otimes |\lambda\rangle\!\langle \lambda|_\Lambda .
\end{equation}
Here, the coefficients $q_{\lambda,n}$ depend only on the spectrum $x$, not on $g$.

We now pass to an easier assisted task.
In the Schur basis, we grant the encoder and decoder free access to the classical register $\Lambda$, and we allow them to discard the multiplicity state $\tau_\lambda$ before compression and recreate it for free after decompression.
Since this only adds side information and free ancillas, the assisted task cannot require more quantum memory than the original one.
Therefore, any lower bound for the assisted task is also a lower bound for the original task.

Fix the same memory $M_n$.
Because the assisted task reveals $\lambda$ and removes the $g$-independent multiplicity state, we may choose the assisted code blockwise:
for each $\lambda$ there is a $\lambda$-visible sector code $(\mathcal E_{\lambda,n},\mathcal D_{\lambda,n})$ acting on the irrep space $\h_\lambda$ and using the same memory $M_n$.
Let
\begin{equation}
    \delta_{\lambda,n}(g):= 
    \frac12\left\|
        \mathcal D_{\lambda,n}\circ \mathcal E_{\lambda,n}(\rho_{g,\lambda})
        -
        \rho_{g,\lambda}
    \right\|_1 .
\end{equation}
Since the assisted task is easier, we can choose these sector codes so that the overall assisted error is at most $\delta_n$.
Using the orthogonality of the Schur blocks in \eqref{eq:converse_schur_form}, this means
\begin{equation}\label{eq:sector_error_average}
    \sup_{g\in \mathrm U(d)}
    \sum_{\lambda:\,\ell(\lambda)\le r}
    q_{\lambda,n}\, \delta_{\lambda,n}(g)
    \le \delta_n .
\end{equation}

Average \eqref{eq:sector_error_average} over the Haar measure on $\mathrm U(d)$ and define
\begin{equation}
    \overline \delta_{\lambda,n}
    := 
    \int_{\mathrm U(d)} \dd g\; \delta_{\lambda,n}(g) .
\end{equation}
Then
\begin{equation}\label{eq:sector_average_bound}
    \sum_{\lambda:\,\ell(\lambda)\le r}
    q_{\lambda,n}\,\overline \delta_{\lambda,n}
    \le \delta_n .
\end{equation}

Define the good set
\begin{equation}
    \mathcal G_n
    := 
    \left\{
        \lambda:\,\ell(\lambda)\le r,\ \overline \delta_{\lambda,n}\le \sqrt{\delta_n}
    \right\}.
\end{equation}
By Markov's inequality and \eqref{eq:sector_average_bound},
\begin{equation}\label{eq:good_mass}
    \sum_{\lambda\notin \mathcal G_n} q_{\lambda,n}
    \le \sqrt{\delta_n},
    \qquad\text{so}\qquad
    \sum_{\lambda\in \mathcal G_n} q_{\lambda,n}
    \ge 1-\sqrt{\delta_n}.
\end{equation}

Next, choose the typical set with $\Delta_n:= \log n$:
\begin{equation}
    \mathcal T_{p,n}:= 
    \left\{
        \lambda:\;
        \lambda_i=0\ \text{for } i>r,\quad
        \max_{1\le i\le r} |\lambda_i-nx_i|
        \le \sqrt n\,\Delta_n
    \right\}.
\end{equation}
If $\lambda\notin \mathcal T_{p,n}$, then for some $1\le i\le r$,
\[
|\lambda_i/n-x_i|>\frac{\Delta_n}{\sqrt n},
\]
and therefore the total variation distance satisfies
\[
d(\lambda/n,p)=
\frac12\sum_{i=1}^d |\lambda_i/n-x_i|
\ge \frac{\Delta_n}{2\sqrt n}.
\]
Applying Lemma~\ref{lem:tail_prob}, we obtain
\begin{equation}\label{eq:typical_mass_bound}
    \sum_{\lambda\notin \mathcal T_{p,n}} q_{\lambda,n}
    \le
    (n+1)^{\frac{r(r+1)}{2}}
    e^{-\Delta_n^2/2}
    = o(1).
\end{equation}

Combining \eqref{eq:good_mass} and \eqref{eq:typical_mass_bound}, we conclude that for all sufficiently large $n$,
\[
\sum_{\lambda\in \mathcal G_n} q_{\lambda,n} +
\sum_{\lambda\in \mathcal T_{p,n}} q_{\lambda,n} >1\ .
\]
Hence, the intersection $\mathcal G_n\cap \mathcal T_{p,n}$ is non-empty. Let $\lambda^{(n)}$ be any sequence of typical partitions in this intersection.

For this specific sequence, the corresponding sector code has Haar-average error
\begin{equation}
\overline \delta_{\lambda^{(n)},n}\le \sqrt{\delta_n}\xrightarrow[n\to\infty]{}0.
\end{equation}

Applying Proposition~\ref{prop:orbit_sector_converse} to the $\lambda$-visible relaxation of that sector task establishes that the memory must be at least the dimension of this specific irrep:
\begin{equation}\label{eq:sector_memory_lower_bound}
    |M_n|
    \ge
    \log \dim \h_{\lambda^{(n)}} + o(1).
\end{equation}

Because $\lambda^{(n)}$ is by definition in the typical set $\mathcal T_{p,n}$, this directly implies the general converse bound for the memory size:
\begin{equation}
    |M_n| \ge \min_{\lambda \in \mathcal T_{p,n}} \log \dim \h_{\lambda} + o(1).
\end{equation}

To establish the specialized asymptotic bound, we now assume the non-zero eigenvalues are strictly decreasing, $x_1>\cdots>x_r>0$.
Because $\lambda^{(n)}\in \mathcal T_{p,n}$, we have
\begin{equation}
\lambda^{(n)}_i = nx_i + O(\sqrt n\,\Delta_n)
\qquad (1\le i\le r),
\qquad
\lambda^{(n)}_i=0
\qquad (i>r).
\end{equation}
Since the $x_i$ are distinct, Lemma~\ref{lem:weyl_asymptotic} applies directly with $\Delta_n=\log n$, yielding the asymptotic expansion of the Weyl dimension formula:
\begin{multline}
    \log \dim \h_{\lambda^{(n)}} =
    \frac{r(2d-r-1)}{2}\log n \\
    + \sum_{1\le i<j\le r}\log (x_i-x_j)
    + (d-r)\sum_{i=1}^r \log x_i
    - \sum_{k=d-r}^{d-1}\log k!
    + O\!\left(\frac{\log n}{\sqrt n}\right).
\end{multline}

Since $\log n/\sqrt n=o(1)$, substituting this expansion into \eqref{eq:sector_memory_lower_bound} yields the final result:
\[
    |M_n|\ge\frac{r(2d-r-1)}{2}\log n
    + \sum_{1\le i<j\le r}\log (x_i-x_j)
    + (d-r)\sum_{i=1}^r \log x_i
    - \sum_{k=d-r}^{d-1}\log k!
    + o(1).
\]
\end{proof}


\section{Generalizing cloning fidelity}\label{sec:fidelity}

In this section, we prove Theorem~\ref{thm:main} that characterizes the cloning fidelity. Let us restate it here.

\begin{mthm}{1}(Cloning Fidelity)
    Let $\rho$ be a density matrix on $\mathbb{C}^d$ with rank $r \le d$. We assume its non-zero eigenvalues are strictly non-degenerate, satisfying $x_1 > x_2 > \dots > x_r > 0$. Let $\mu, \nu$ be two ordered partitions labeling $\mathrm{GL}(d,\mathbb C)$ irreps such that $(\nu-\mu)$ is a dominant weight. Their sizes are parameterized by a large integer $n$ such that $|\mu|,|\nu|=\Theta(n)$. Let the distance between the irreps be $d(\mu,\nu) := \sum_{i=1}^d |\mu_i-\nu_i|$. Let $\rho_{\mu}$ and $\rho_\nu$ be the normalized $\mathrm{GL}(d,\mathbb C)$ representations of $\rho$. Let $\C_{\mu\to\nu}$ be the generalized cloning map from $\mu$ to $\nu$. For any arbitrarily small $\epsilon>0$, the cloning fidelity is asymptotically bounded by:
    \begin{equation}
        F(\C_{\mu\to\nu}(\rho_\mu), \rho_\nu)\ge 1-O\left(\frac{d(\mu,\nu)}{n^{1-\epsilon}}\right)\ ,\quad F(\C_{\nu\to\mu}(\rho_\nu), \rho_\mu)\ge 1-O\left(\frac{d(\mu,\nu)}{n^{1-\epsilon}}\right) \ .
    \end{equation}
\end{mthm}

We sketch the proof outline of Theorem~\ref{thm:main}. Casting the channel $\C_{\mu\to\nu}$ in the Kraus representation, we observe that each Kraus operator corresponds to a vector $\ket{W}$ in the irrep $\omega$. We simplify the problem by only estimating the contribution of the Kraus operator corresponding to the highest weight vector $\ket{\omega}$ in $\omega$. With the help of Lemma~\ref{lem:monotonicity}, we establish that this gives a lower bound to the generalized cloning fidelity.

In evaluating the simplified fidelity, the contribution attributes a quantum part and a classical part. The classical contribution corresponds to the ratio of the spectra of $\rho_\mu$ and $\rho_\nu$, and it is estimated in Lemma~\ref{lem:ratio}. The quantum contribution is determined how to the eigenspaces (weight spaces) of $\rho_\mu$ and $\rho_\nu$ align with each other. We show in Lemma~\ref{lem:perturbation} that the mismatch between the eigenspaces is vanishingly small.

Furthermore, since most of the spectral weight of $\rho_\mu$ and $\rho_\nu$ is concentrated on height weights, we truncate the lower weights to simplify the analysis. The truncation error is estimated in Lemma~\ref{lem:tail}. 

Finally, we put all of this together to prove the full Theorem~\ref{thm:main}. Before doing the proof, we need to go over some technical preliminaries\\

\subsection{Technical Preliminaries}\label{sec:prelim}

We review the representation theory of $\mathrm{GL}(d,\mathbb C)$ and its complexified Lie algebra $\mathfrak{gl}(d, \mathbb{C})$, and other relevant tools need for our proofs. We refer to a standard text~\cite{fulton1991representation} for more background on representation theory.

\begin{itemize}
    
\item\textbf{Roots, Coroots, and the Cartan-Weyl Basis.}
We fix a standard Cartan subalgebra $\mathfrak{h}$ of diagonal matrices, denoting the dual weight space by $\mathfrak{h}^*$, equipped with the standard Euclidean inner product $(\cdot, \cdot)$ induced by the trace.
We choose the standard Cartan-Weyl basis $\{H_i, E_\alpha, E_{-\alpha}\}$. In the canonical basis $\{e_i\}_{i=1}^d$ of $\mathbb{R}^d$, the positive roots are $\alpha = e_i - e_j$ for $i < j$, forming the set $\Phi^+$. The simple roots are $\alpha_i = e_i - e_{i+1}$ for $i = 1, \dots, d-1$, and they generate the positive root lattice $Q_+$. Because $\mathfrak{gl}(d, \mathbb{C})$ belongs to the type $A_{d-1}$ Lie algebra family, roots and coroots are canonically identified ($\alpha^\vee = \alpha$). For simple roots, the pairing with any weight $\lambda \in \mathfrak{h}^*$ extracts the difference between adjacent Dynkin labels:
\begin{equation}\label{eq:pairing}
    \langle \lambda, \alpha_i^\vee \rangle = (\lambda, \alpha_i) = \lambda_i - \lambda_{i+1} \ .
\end{equation}
For a compound root \(\eta=\sum_k c_k\alpha_k\), we define the pairing of a weight with a compound coroot to be the weighted sum of its pairings with the simple coroots: \[ \langle \lambda,\eta^\vee\rangle := \sum_k c_k\langle \lambda,\alpha_k^\vee\rangle \ . \]

\item\textbf{The Quadratic Casimir Operator.}
The total quadratic Casimir operator $C_2$ acts as a scalar on any irrep $\h_\lambda$. In the Cartan-Weyl basis, 
\begin{equation}
    C_2 = \sum_{i=1}^d H_i^2 + \sum_{\alpha\in\Phi^+} (E_{-\alpha} E_\alpha + E_\alpha E_{-\alpha}) \ .
\end{equation}
The exact eigenvalue on $\h_\lambda$ is $C_2(\lambda) = (\lambda, \lambda + 2\varrho)$, where $\varrho = \frac{1}{2} \sum_{\alpha \in\Phi^+} \alpha$ is the Weyl vector.


\item\textbf{Weight Spaces, Root Paths, and Height.}
Let $\h_\mu$ denote the finite-dimensional irrep of $\mathrm{GL}(d,\mathbb{C})$ labeled by its strictly dominant highest weight $\mu$. The space decomposes into orthogonal weight spaces $W_\mu(\delta)$ of weight $\mu-\delta$, where the weight offset $\delta \in Q_+$ can be expanded in the simple root basis as $\delta = \sum_{k=1}^{d-1} c_k \alpha_k$ with non-negative integers~$c_k$.

We define the \emph{height} of this weight offset as the sum of its simple root coefficients:
\begin{equation}
    \mathrm{ht}(\delta) := \sum_{k=1}^{d-1} c_k \ .
\end{equation}
We can equivalently describe this offset via a sequence, or \emph{path}, of simple roots descending from the highest weight. The total number of steps in such a path exactly matches the height. Therefore, we will frequently use $|\delta|$ to denote both the path length and the height $\mathrm{ht}(\delta)$ interchangeably.

Rather than relying on an explicit combinatorial construction for the states, we simply denote an arbitrary orthonormal basis for a given weight space $W_\mu(\delta)$ as $\{\ket{\mu, \delta, k}\}_{k=1}^{m_\mu(\delta)}$. Here, the index $k$ resolves the inner degeneracy of the weight space, which is given exactly by the Kostka number: $m_\mu(\delta) = K_{\mu,\,\mu-\delta}$.

\item\textbf{\(\alpha\)-strings and local \(\mathfrak{sl}_2\) structure.}
Fix a positive root \(\alpha\in \Phi^+\). The operators
\[
E_\alpha,\qquad E_{-\alpha},\qquad H_\alpha:=[E_\alpha,E_{-\alpha}]
\]
span an \(\mathfrak{sl}_2\)-subalgebra. Accordingly, when a finite-dimensional irrep \(\h_\lambda\) of \(\mathrm{GL}(d,\mathbb C)\) is restricted to this rank-one subalgebra, it decomposes as a direct sum of irreducible \(\mathfrak{sl}_2\)-modules. We call each such summand an \emph{\(\alpha\)-string}.

Let \(\ket{v_{\mathrm{top}}}\) be a highest-weight vector in one \(\alpha\)-string, so that \(E_\alpha \ket{v_{\mathrm{top}}}=0\), and suppose \(\ket{v_{\mathrm{top}}}\) has weight \(\beta\). The local highest weight of this \(\mathfrak{sl}_2\)-module is
\begin{equation}
    m:=\langle \beta,\alpha^\vee\rangle \in \mathbb Z_{\ge 0}.
\end{equation}
Some authors would call \(m+1\) the length of the string; we will use \(m\) as the \emph{length parameter}. The weights in the string are
\[
\beta,\ \beta-\alpha,\ \beta-2\alpha,\ \dots,\ \beta-m\alpha.
\]

Choose an orthonormal basis \(\{\ket{v_q}\}_{q=0}^m\) of this string, where \(\ket{v_q}\) has weight \(\beta-q\alpha\). The standard \(\mathfrak{sl}_2\) matrix-element formulas are
\begin{equation}\label{eq:alpha_string_minus}
    E_{-\alpha}\ket{v_q} = c_q^-\, \ket{v_{q+1}},
    \qquad
    |c_q^-|^2=(q+1)(m-q),
\end{equation}
and
\begin{equation}\label{eq:alpha_string_plus}
    E_{\alpha}\ket{v_q} = c_q^+\, \ket{v_{q-1}},
    \qquad
    |c_q^+|^2=q(m-q+1),
\end{equation}
with the conventions \(\ket{v_{-1}}=\ket{v_{m+1}}=0\).

The integer \(q\) measures how far the vector lies below the local top of the \(\alpha\)-string. If one restricts to a subspace spanned by weights of depth at most \(T\) below the global highest weight, then necessarily
\[
q\le T,
\]
because each application of \(E_{-\alpha}\) lowers the weight by \(\alpha\), whose height satisfies \(|\alpha|\ge 1\). Therefore, on such a restricted subspace,
\begin{equation}\label{eq:alpha_string_norm_bound}
    \max_{0\le q\le \min\{T,m\}}\{|c_q^-|,|c_q^+|\}
    = O(\sqrt{Tm}).
\end{equation}

For type \(A_{d-1}\), every positive root has the form \(\alpha=e_i-e_j\) with \(i<j\). If \(\beta\) is the local top weight of an \(\alpha\)-string, then
\[
m=\langle \beta,\alpha^\vee\rangle=\beta_i-\beta_j\le \max_{a,b}(\beta_a-\beta_b).
\]
Since every weight \(\beta\) of \(\h_\lambda\) lies in the convex hull of the Weyl orbit \(W\lambda\), each coordinate of \(\beta\) lies between \(\lambda_d\) and \(\lambda_1\). Hence
\begin{equation}\label{eq:alpha_string_length_bound}
    m\le \lambda_1-\lambda_d.
\end{equation}

\item\textbf{Verma Modules and Kostant's Partition Function.}
Let $U(\mathfrak{gl}(d,\mathbb C))$ be the universal enveloping algebra of $\mathfrak{gl}(d,\mathbb C)$. For a dominant integral weight $\lambda$, the Verma module $M(\lambda)$ is the infinite-dimensional highest-weight cyclic module generated by a highest-weight vector $\ket{\lambda}$ satisfying $E_\alpha \ket{\lambda} = 0$ for all $\alpha \in \Phi^+$ and $H_i \ket{\lambda} = \lambda_i \ket{\lambda}$. 

The weight spaces of $M(\lambda)$ correspond to weights of the form $\lambda - \delta$ for $\delta \in Q_+$. The dimension of each weight space is strictly independent of $\lambda$ and is given exactly by Kostant's partition function $K(\delta)$, which counts the number of ways to express the offset $\delta$ as a sum of positive roots.

The finite-dimensional irrep $\h_\lambda$ is the unique irreducible quotient of the Verma module $M(\lambda)$. The kernel of the projection map $M(\lambda) \to \h_\lambda$ is generated by singular vectors that appear at depth $\langle \lambda, \alpha_i^\vee \rangle + 1$ for each simple root $\alpha_i$. Therefore, for shallow weight offsets satisfying $|\delta| \le \min_i \langle \lambda, \alpha_i^\vee \rangle$, no states are quotiented out. In this regime, the weight spaces of $\h_\lambda$ are perfectly isomorphic to those of $M(\lambda)$, yielding an inner degeneracy exactly equal to $K(\delta)$. We rely on this equivalence for the multiplicities in the shallow-depth regime.

\begin{lem}[Kostka Numbers]\label{lem:kostka}
    Let $\mu,\omega$ be dominant integral highest weights of $\mathrm{GL}(d,\mathbb C)$, let $\delta\in Q_+$ be a weight offset, and let $g_\mu:= \min_{1\le i\le d-1}(\mu_i-\mu_{i+1})$ be the edge gap of $\mu$. Then the inner degeneracies (Kostka numbers) satisfy
    \begin{equation}
        K_{\mu,\mu-\delta} \le K_{\mu+\omega,\mu+\omega-\delta}.
    \end{equation}
    Moreover, equality holds when $|\delta|\le g_\mu$.
\end{lem}
\begin{proof}See Appendix~\ref{app:proofs}. \end{proof}

\item\textbf{$\mathrm{GL}(d,\mathbb C)$ irreps of density matrices.} 
The normalized irrep of $\rho$, defined as $\rho_\mu := \pi_\mu(\rho)/s_\mu(x)$, is block-diagonalized in weights: 
\begin{equation}
    \rho_\mu = \sum_{\delta\in Q_+} p_\mu(\delta) \Pi_{\mu-\delta} \ .
\end{equation}
where \(Q_+\) is the positive root cone of \(\mathrm{GL}(r)\), \(\Pi_{\mu-\delta}\) projects onto the weight space of weight \(\mu-\delta\). 

The eigenvalues are $p_\mu(\delta) := x^{\mu-\delta}/s_\mu(x)$, where $x=(x_i)_{i=1}^r$ denotes the non-zero spectrum of $\rho$, and $x^w := \prod_{i=1}^r x_i^{w_i}$. Note that $p_\mu(\delta)$ is not a probability distribution over $\delta$, but with multiplicities included,  $p_\mu(\delta)m_\mu(\delta)$  is a probability distribution.

The normalization factor $s_{\mu}(x) := \Tr\, \pi_{\mu}(\rho)$ is the Schur polynomial. We utilize both its combinatorial expansion over symmetric groups and its classical determinantal formula over the $r$ non-zero variables:
\begin{equation}
    s_\mu(x) = \frac{\sum_{\pi\in S_d} \mathrm{sgn}(\pi)\ \prod_{i=1}^d x_i^{\mu_{\pi(i)}-\delta_{\pi(i)}+d-\pi(i)}}{\prod_{i<j}(x_i-x_j)} = \frac{\det(x_i^{\mu_j+r-j})}{\det(x_i^{r-j})} \ .
\end{equation}
For more on Schur polynomials, see~\cite{macdonald1995symmetric}.



\item\textbf{Asymptotic Parameters.}
Let the width of a Young diagram $\lambda$ be $\|\lambda\|:=\lambda_1-\lambda_d$ Because $\rho$ has rank $r \le d$, we restrict to the subgroup $\mathrm{GL}(r)$ and define the \emph{edge gap} as:
\begin{equation}
    g_\lambda := \min_{1\leq i\leq r-1}\langle \lambda, \alpha_i^\vee \rangle= \min_{1\leq i\leq r-1}(\lambda_i-\lambda_{i+1}) \ .
\end{equation}
We operate in the scaling regime where strongly regular highest weights have edge gaps scaling as $g_\mu \approx g_\nu = \Theta(n)$. We also mostly consider the case of \emph{strict dominance}: $\nu-\mu=:\omega$ is also an ordered partition (Young diagram), and $\omega$ is much smaller: $|\omega|, \|\omega\|=o(n)$.


\item\textbf{Principal Angles.}
Let $P$ and $Q$ be orthogonal projectors onto two subspaces $U,V \subseteq \mathbb{C}^n$. The \emph{principal angles} $\theta_1,\dots,\theta_r \in [0,\pi/2]$ between $U$ and $V$ are defined recursively, or equivalently via the singular values of the overlap map between the two subspaces. If $r=\min(\dim U,\dim V)$, then the numbers $\cos\theta_i$ are precisely the singular values associated with the pair $(U,V)$. A convenient operator-theoretic way to encode these angles is through the compression $PQP$ acting on $\operatorname{Im}(P)$: its nonzero eigenvalues are exactly $\cos^2\theta_i$.

\begin{lem}[Principle Angles]\label{lem:principle_angles}
Let $P$ and $Q$ be orthogonal projectors, and let $\theta_1,\dots,\theta_r$ denote the principal angles between $\operatorname{Im}(P)$ and $\operatorname{Im}(Q)$, where $r=\min(\operatorname{rank}P,\operatorname{rank}Q)$. Then
\begin{equation}
    \Tr \sqrt{PQP} = \sum_{i=1}^r \cos\theta_i \ .
\end{equation}
\end{lem}
\begin{proof}
    See Appendix~\ref{app:proofs}.
\end{proof}

\end{itemize}



\subsection{Proof}\label{sec:proof}

\emph{Proof of Theorem~\ref{thm:main}.}
Since \((\nu-\mu)\) is a dominant weight, we have
\[
\omega:= (\nu-\mu)^+=\nu-\mu.
\]
Then the distance is exactly the size of the ancilla offset:
\[
d(\mu,\nu)=|\omega|:=\sum_i\omega_i.
\]
We first bound the \emph{forward-direction} fidelity \(F(\C_{\mu\to\nu}(\rho_\mu),\rho_\nu)\).

Because the perturbative analysis in Lemma~\ref{lem:perturbation} takes place on the product space
\(\mathcal H_\mu\otimes \mathcal H_\omega\), we embed the intrinsic channel
\[
\C_{\mu\to\nu}:\mathcal B(\mathcal H_\mu)\to \mathcal B(\mathcal H_\nu),
\qquad
\C_{\mu\to\nu}(X)=\frac{d_\mu}{d_\nu}V^\dagger(X\otimes \id_\omega)V,
\]
into the larger space using the intertwining isometry \(V:\mathcal H_\nu\to \mathcal H_\mu\otimes \mathcal H_\omega\). Define
\begin{equation}\label{eq:embedded_channel}
    \widetilde{\C}_{\mu\to\nu}(X)
    :=
    V\C_{\mu\to\nu}(X)V^\dagger
    =
    \frac{d_\mu}{d_\nu}P_\nu(X\otimes \id_\omega)P_\nu ,
\end{equation}
where \(P_\nu=VV^\dagger\in\mathcal B(\h_\mu\otimes\h_\omega)\) is the orthogonal projector onto the embedded copy of \(\mathcal H_\nu\).
Likewise, define the embedded target state and embedded weight projectors by
\begin{equation}\label{eq:embedded_target}
    \widetilde{\rho}_\nu
    :=
    V\rho_\nu V^\dagger
    =
    \sum_{\delta\in Q_+} p_\nu(\delta)\,\widetilde{\Pi}_{\nu-\delta},
    \qquad
    \widetilde{\Pi}_{\nu-\delta}
    :=
    V\Pi_{\nu-\delta}V^\dagger .
\end{equation}
Since fidelity is invariant under a common isometric embedding,
\begin{equation}\label{eq:isometric_invariance_forward}
    F(\C_{\mu\to\nu}(\rho_\mu),\rho_\nu)
    =
    F(\widetilde{\C}_{\mu\to\nu}(\rho_\mu),\widetilde{\rho}_\nu).
\end{equation}
From this point on, all forward-direction fidelities are computed on
\(\mathcal H_\mu\otimes \mathcal H_\omega\).

Choose an orthonormal basis \(\{\ket{a}\}\) of \(\mathcal H_\omega\) containing the highest-weight vector \(\ket{\omega}\). Then \(\widetilde{\C}_{\mu\to\nu}\) admits the decomposition
\[
\widetilde{\C}_{\mu\to\nu}(X)=\sum_a \mathcal K_a(X),
\qquad
\mathcal K_a(X):=\frac{d_\mu}{d_\nu}P_\nu(X\otimes \ketbra{a}{a})P_\nu .
\]
In particular, the highest-weight Kraus branch is
\begin{equation}\label{eq:kraus_proj}
    \mathcal K_\omega(X)
    :=
    \frac{d_\mu}{d_\nu}P_\nu(X\otimes \ketbra{\omega}{\omega})P_\nu .
\end{equation}

We now use \(F(\mathcal K_\omega(\rho_\mu),\widetilde{\rho}_\nu)\) as a lower bound on
\(F(\widetilde{\C}_{\mu\to\nu}(\rho_\mu),\widetilde{\rho}_\nu)\). The required monotonicity statement is the following.

\begin{lem}[Monotonicity]\label{lem:monotonicity}
    Let $\mathcal{N}$ be a quantum channel and $\rho = \mathcal{N}(\eta)$ for some input state $\eta \geq 0$. Let $K$ be a Kraus operator of $\mathcal{N}$ such that $\mathcal{N}(\cdot) = K(\cdot)K^\dagger + \sum_{i} L_i (\cdot) L_i^\dagger$, and define the unnormalized state $\rho' = K \eta K^\dagger$. For any density operator $\sigma$, the fidelity satisfies:
    \begin{equation}
        F(\rho', \sigma) \leq F(\rho, \sigma) \ .
    \end{equation}
\end{lem}
\begin{proof}
See Appendix~\ref{app:proofs}.
\end{proof}

%Since \([\C_{\mu\to\nu}(\rho_\mu),\rho_\nu]=0\) by Lemma~\ref{lem:commutativity}, conjugating by \(V(\cdot )V^\dagger\) gives \[ [\widetilde{\C}_{\mu\to\nu}(\rho_\mu),\widetilde{\rho}_\nu]=0\ . \]
Applying Lemma~\ref{lem:monotonicity} to \(\widetilde{\C}_{\mu\to\nu}\), its Kraus branch \(\mathcal K_\omega\), and the commuting pair
\(\widetilde{\C}_{\mu\to\nu}(\rho_\mu),\widetilde{\rho}_\nu\), we obtain
\begin{equation}\label{eq:forward_monotonicity}
    F(\C_{\mu\to\nu}(\rho_\mu),\rho_\nu)
    =
    F(\widetilde{\C}_{\mu\to\nu}(\rho_\mu),\widetilde{\rho}_\nu)
    \ge
    F(\mathcal K_\omega(\rho_\mu),\widetilde{\rho}_\nu).
\end{equation}

Because \(\mathcal K_\omega(\Pi_{\mu-\delta})\) is supported on the total-weight sector \(\nu-\delta\), the fidelity decomposes blockwise:
\begin{equation}\label{eq:block_fidelity}
    F(\C_{\mu\to\nu}(\rho_\mu),\rho_\nu)
    \ge
    \sum_{\delta\in Q_+}
    \sqrt{p_\mu(\delta)p_\nu(\delta)}\,
    F(\mathcal K_\omega(\Pi_{\mu-\delta}),\widetilde{\Pi}_{\nu-\delta}) .
\end{equation}
Thus the problem splits into a quantum part, namely the block fidelity
\(F(\mathcal K_\omega(\Pi_{\mu-\delta}),\widetilde{\Pi}_{\nu-\delta})\), and a classical part, namely the Bhattacharyya coefficient:\footnote{Neither $p_\mu(\delta)$ nor $p_\nu(\delta)$ is a probability distribution, so we are slightly abusing the term here.} \(\sum_\delta \sqrt{p_\mu(\delta)p_\nu(\delta)}\).\\

We proceed by first estimating the block-wise quantum fidelity in~\eqref{eq:block_fidelity}. The key tool is the Davis--Kahan \(\sin\theta\) theorem.

\begin{lem}[Davis-Kahan Theorem~\cite{davis1970rotation}]\label{lem:davis_kahan}
    Let \(H_0\) and \(H = H_0 + V\) be Hermitian operators on a finite-dimensional Hilbert space \(\mathcal{H}\). Let the spectrum of \(H_0\) be partitioned into two disjoint sets \(S_0\) and \(S_1\), separated by a spectral gap \(\Delta := \min_{\lambda \in S_0,\, \lambda' \in S_1} |\lambda - \lambda'|\).
    Let \(P_0\) be the orthogonal projector onto the eigenspace of \(H_0\) corresponding to \(S_0\).
    If \(\|V\| < \Delta\), then the eigenvalues of \(H\) split into two sets \(\sigma_0(H)\) and \(\sigma_1(H)\) of the same cardinality as \(S_0\) and \(S_1\) respectively, separated by a gap of at least \(\Delta - \|V\|\).
    Let \(P\) be the orthogonal projector onto the eigenspace of \(H\) corresponding to \(\sigma_0(H)\). Then
    \begin{equation}
        \sin\theta_{\max}\equiv \|P - P_0\|\le \frac{\|V P_0\|}{\Delta - \|V\|} \,,
    \end{equation}
    where \(\theta_{\max}\) is the maximum principal angle between the subspaces \(\mathrm{im}(P)\) and \(\mathrm{im}(P_0)\).
\end{lem}

\begin{rem}
    The projector \(P\) can be defined explicitly via the resolvent $(z-H)^{-1}$ and Cauchy integral. Let \(\Gamma\) be a positively oriented contour in the complex plane that encloses \(S_0\) while excluding \(S_1\). The condition \(\|V\| < \Delta\) ensures that no eigenvalue of \(H\) crosses \(\Gamma\), so that we can define $P$ using the resolvent
    \begin{equation}
        P = -\frac{1}{2\pi i}\oint_\Gamma (z - H)^{-1}\,dz
    \end{equation}
    and it has the same rank as \(P_0\). This definition makes explicit that \(P\) depends continuously on \(V\) via the resolvent and reduces to \(P_0= -\frac{1}{2\pi i}\oint_\Gamma (z - H_0)^{-1}\,dz\) when \(V = 0\).
\end{rem}

We use the Davis-Kahan Theorem to prove a key technical lemma that controls the quantum contribution to the fidelity.

\begin{lem}[Highest-weight subspace perturbation]\label{lem:perturbation}
Let \(\nu\), \(\mu\), and \(\omega\) be dominant highest weights of \(\mathrm{GL}(d,\mathbb C)\) satisfying the strict dominance relation $\nu=\mu+\omega$.
Let \(P_\nu\) be the orthogonal projector onto the copy of \(\mathcal H_\nu\subset \mathcal H_\mu\otimes\mathcal H_\omega\), and let
\[
P_\mu^{(\omega)}:=\id_\mu\otimes \ketbra{\omega}{\omega},
\]
where \(\ket{\omega}\in\mathcal H_\omega\) is a highest-weight vector.
For a positive-root sum \(\delta\in Q_+\) of height \(|\delta|\), let \(Q_{\nu-\delta}\) denote the orthogonal projector onto the total-weight subspace $W_{\nu-\delta}\subset \mathcal H_\mu\otimes \mathcal H_\omega$.
Since both \(P_\nu\) and \(P_\mu^{(\omega)}\) commute with the total Cartan action, they commute with \(Q_{\nu-\delta}\). We therefore define
\begin{equation}
    \widetilde{\Pi}_{\nu-\delta}:=Q_{\nu-\delta}P_\nu=P_\nu Q_{\nu-\delta},
    \qquad
    \Pi_{\mu-\delta}^{(\omega)}:=Q_{\nu-\delta}P_\mu^{(\omega)}=P_\mu^{(\omega)}Q_{\nu-\delta}.
\end{equation}
Then
\begin{equation}
    \left\|\widetilde{\Pi}_{\nu-\delta}-\Pi_{\mu-\delta}^{(\omega)}\right\|
    = O\!\left(\sqrt{\frac{|\delta| \|\omega\|}{n}}\right)
\end{equation}
under the scaling regime
\[
|\mu|=\Theta(n),\qquad g_\mu=\Theta(n),\qquad |\delta|^2\|\omega\|\ll n,
\]
where \(\|\omega\|:=\omega_1-\omega_d\) denotes the width of \(\omega\).
\end{lem}

\begin{rem}
The condition \(|\delta|^2\|\omega\|\ll n\) implies \(|\delta|=o(n)\). Since \(g_\mu=\Theta(n)\), for all sufficiently large \(n\) one has \(|\delta|\le g_\mu\), and Lemma~\ref{lem:kostka} yields
\[
K_{\mu,\mu-\delta}=K_{\nu,\nu-\delta}.
\]
Hence the two restricted projectors have the same rank, as they should.
\end{rem}

\begin{proof}
We work entirely on the fixed total-weight slice
\[
W_{\nu-\delta}\subset \mathcal H_\mu\otimes \mathcal H_\omega
\]
and apply the Davis--Kahan theorem there. The idea is to view the total quadratic Casimir of $\mu\otimes\omega$ as an interacting Hamiltonian
\[
H=H_0+V,
\]
where \(H_0\) is the diagonal Cartan part and \(V\) contains the root-exchange terms. The distinguished eigenspace of \(H_0\) will be the uncoupled subspace \(\Pi_{\mu-\delta}^{(\omega)}\), while the corresponding spectral projector of the perturbed \(H\) will be the exact coupled subspace \(\widetilde{\Pi}_{\nu-\delta}\).

\medskip
\noindent\textbf{Step 1: Restrict the total Casimir and identify the unperturbed eigenspace.}

For notational clarity, we let
\[
Q:=Q_{\nu-\delta}.
\]
Write the total quadratic Casimir on \(\mathcal H_\mu\otimes\mathcal H_\omega\) in the form
\begin{equation}
\begin{aligned}
    C_2(\mu\otimes\omega)
    &=
    C_2(\mu)\,\id_\mu\otimes \id_\omega
    + \id_\mu\otimes C_2(\omega)\,\id_\omega
    + 2\sum_{i=1}^d H_i^{(\mu)}\otimes H_i^{(\omega)} \\
    &\quad
    + 2\sum_{\alpha>0}
    \Bigl(
    E_{-\alpha}^{(\mu)}\otimes E_\alpha^{(\omega)}
    +
    E_\alpha^{(\mu)}\otimes E_{-\alpha}^{(\omega)}
    \Bigr).
\end{aligned}
\end{equation}
Since \(C_2(\mu\otimes\omega)\) commutes with the total Cartan generators, it preserves every total-weight space. Thus the operator relevant for the slice \(W_{\nu-\delta}\) is
\[
H:=Q\,C_2(\mu\otimes\omega)\,Q=C_2(\mu\otimes\omega)\big|_{W_{\nu-\delta}}.
\]
We split it as \(H=H_0+V\), where
\begin{equation}
    H_0
    :=
    Q\Bigl(
    C_2(\mu)\,\id_\mu\otimes \id_\omega
    + \id_\mu\otimes C_2(\omega)\,\id_\omega
    + 2\sum_{i=1}^d H_i^{(\mu)}\otimes H_i^{(\omega)}
    \Bigr)Q,
\end{equation}
and
\begin{equation}
    V
    :=
    Q\Bigl(
    2\sum_{\alpha>0}
    \bigl(
    E_{-\alpha}^{(\mu)}\otimes E_\alpha^{(\omega)}
    +
    E_\alpha^{(\mu)}\otimes E_{-\alpha}^{(\omega)}
    \bigr)
    \Bigr)Q.
\end{equation}

Now consider the subspace defined by the projector
\[
P_0:=\Pi_{\mu-\delta}^{(\omega)}.
\]
Because \(\ket{\omega}\) has weight \(\omega\), fixing the total weight to be \(\nu-\delta=\mu+\omega-\delta\) forces the \(\mu\)-register to have weight \(\mu-\delta\). Therefore
\begin{equation}\label{eq:proj_mu_omega}
    \Pi_{\mu-\delta}^{(\omega)}
    =
    \Pi_{\mu-\delta}\otimes \ketbra{\omega}{\omega}.
\end{equation}
On this subspace, \(H_0\) acts as the scalar
\begin{equation}\label{eq:lambda0_def}
    \lambda_0
    =
    C_2(\mu)+C_2(\omega)+2(\mu-\delta,\omega).
\end{equation}

\medskip
\noindent\textbf{Step 2: The unperturbed spectral gap is of order \(n\).}

Let \(v\in W_{\nu-\delta}\) be orthogonal to \(\mathrm{im}(P_0)\). Then the \(\omega\)-register must have some lower weight
\[
\omega-\eta,
\qquad
\eta\in Q_+,\ \eta\neq 0.
\]
Since the total weight is fixed to \(\nu-\delta\), the \(\mu\)-register must then have weight
\[
\mu-\delta+\eta.
\]
Because this must be an actual weight of \(\mathcal H_\mu\), one necessarily has
\[
\delta-\eta\in Q_+,
\qquad\text{hence}\qquad
\eta\preceq \delta
\quad\text{and}\quad
|\eta|\le |\delta|.
\]

On such a product weight space, \(H_0\) acts diagonally with eigenvalue
\begin{equation}
    \lambda_1(\eta)
    =
    C_2(\mu)+C_2(\omega)+2(\mu-\delta+\eta,\omega-\eta).
\end{equation}
Subtracting from \eqref{eq:lambda0_def}, we get
\begin{align}
    \lambda_0-\lambda_1(\eta)
    &=
    2(\mu-\delta,\omega)-2(\mu-\delta+\eta,\omega-\eta)\notag\\
    &=
    2(\mu,\eta)-2(\delta,\eta)-2(\omega,\eta)+2(\eta,\eta).
    \label{eq:gap_difference}
\end{align}

We now estimate the terms on the right-hand side.

Write
\[
\eta=\sum_k c_k\alpha_k,
\qquad
\delta=\sum_i b_i\alpha_i,
\qquad
c_k,b_i\in\mathbb Z_{\ge 0}.
\]
Then
\begin{equation}
    (\mu,\eta)
    =
    \sum_k c_k(\mu,\alpha_k)
    =
    \sum_k c_k(\mu_k-\mu_{k+1})
    \ge
    g_\mu \sum_k c_k
    =
    g_\mu |\eta|.
\end{equation}
Next,
\begin{equation}
\begin{aligned}
    |(\delta,\eta)|
    &=
    \left|
    \sum_{i,k} b_i c_k (\alpha_i,\alpha_k)
    \right|
    \le
    \sum_{i,k} b_i c_k |(\alpha_i,\alpha_k)| \\
    &\le
    2\sum_{i,k} b_i c_k
    =
    2|\delta|\,|\eta|,
\end{aligned}
\end{equation}
since in type \(A_{d-1}\) one has \((\alpha_i,\alpha_k)\in\{-1,0,2\}\).
Finally,
\begin{equation}
\begin{aligned}
    (\omega,\eta)
    &=\sum_k c_k(\omega,\alpha_k)
    =\sum_k c_k(\omega_k-\omega_{k+1}) \\
    &\le \Bigl(\omega_1-\omega_d\Bigr)\sum_k c_k
    = \|\omega\| |\eta|.
\end{aligned}
\end{equation}
Substituting these bounds into \eqref{eq:gap_difference} and using \((\eta,\eta)\ge 0\), we obtain
\begin{equation}\label{eq:gap_lower_bound_intermediate}
    \lambda_0-\lambda_1(\eta)
    \ge
    |\eta|\Bigl(2g_\mu-4|\delta|-2\|\omega\|\Bigr).
\end{equation}

Now the nontrivial regime \(|\delta|\ge 1\), \(\|\omega\|>0\), together with \(|\delta|^2\|\omega\|\ll n\), implies both
\[
|\delta|=o(n),
\qquad
\|\omega\|=o(n).
\]
Since \(g_\mu=\Theta(n)\), the bracket in \eqref{eq:gap_lower_bound_intermediate} is positive and of order \(n\) for all sufficiently large \(n\). Therefore the spectral gap of \(H_0\) between the eigenspace \(P_0\) and the rest of the slice is
\begin{equation}\label{eq:Delta_def}
    \Delta
    :=
    \mathrm{dist}\Bigl(\{\lambda_0\},\mathrm{spec}(H_0)\setminus\{\lambda_0\}\Bigr)
    =\min_{\eta\neq 0}(\lambda_0-\lambda_1(\eta)) 
    = \Omega(n).
\end{equation}

\medskip
\noindent\textbf{Step 3: The corresponding perturbed projector is \(\widetilde{\Pi}_{\nu-\delta}\).}

Let \(P\) denote the spectral projector of \(H\) onto the interval
\[
I:=\bigl[\lambda_0-\Delta/2,\ \lambda_0+\Delta/2\bigr].
\]
We claim that
\begin{equation}\label{eq:P_equals_restricted_nu}
    P=\widetilde{\Pi}_{\nu-\delta}=Q_{\nu-\delta}P_\nu.
\end{equation}

Indeed, because \(H\) is the restriction of the total Casimir to the slice, its eigenspaces are precisely
\[
Q_{\nu-\delta}P_\lambda
\]
for those irreps \(\lambda\subset \mu\otimes\omega\) that have a nonzero \((\nu-\delta)\)-weight space. In particular, \(Q_{\nu-\delta}P_\nu\) is the eigenspace of \(H\) with eigenvalue \(C_2(\nu)\).

First, \(C_2(\nu)\) lies in \(I\). Using \(C_2(\nu)=(\nu,\nu+2\varrho)\), where \(\varrho\) is the Weyl vector, and \(\nu=\mu+\omega\), one checks directly that
\begin{equation}
    \lambda_0-C_2(\nu)
    =
    -2(\delta,\omega).
\end{equation}
Hence
\begin{equation}
    |\lambda_0-C_2(\nu)|
    \le
    2|\delta| \|\omega\|
    =
    o(n),
\end{equation}
so \(C_2(\nu)\in I\) for all sufficiently large \(n\).

Next, no other irrep contributes an eigenvalue in \(I\). Suppose \(\nu-\gamma\) is another highest weight such that
\[
Q_{\nu-\delta}P_{\nu-\gamma}\neq 0.
\]
Then \(\nu-\delta\) must be a weight of the irrep \(\nu-\gamma\), so
\[
(\nu-\gamma)-(\nu-\delta)=\delta-\gamma\in Q_+.
\]
Hence \(\gamma\preceq \delta\), and therefore \(|\gamma|\le |\delta|\).

For \(\gamma\neq 0\),
\begin{align}
    C_2(\nu)-C_2(\nu-\gamma)
    &=
    2(\nu+\varrho,\gamma)-(\gamma,\gamma)\notag\\
    &=
    2(\nu,\gamma)+2(\varrho,\gamma)-(\gamma,\gamma).
\end{align}
Because \(\nu=\mu+\omega\) is dominant and \(g_\nu\ge g_\mu\), we have
\[
(\nu,\gamma)\ge g_\nu |\gamma|\ge g_\mu |\gamma|.
\]
Moreover, since \(\varrho\) depends only on \(d\), one has \((\varrho,\gamma)=O(|\gamma|)\), and also \((\gamma,\gamma)=O(|\gamma|^2)\). Therefore
\begin{equation}
    C_2(\nu)-C_2(\nu-\gamma)
    \ge
    2g_\mu |\gamma| - O(|\gamma|)-O(|\gamma|^2).
\end{equation}
Since \(g_\mu=\Theta(n)\) and \(|\gamma|\le |\delta|=o(n)\), the right-hand side is \(\Omega(n)\) for every \(\gamma\neq 0\). Thus every competing eigenvalue lies a distance of order \(n\) below \(C_2(\nu)\), and hence outside \(I\).

This proves \eqref{eq:P_equals_restricted_nu}.

\medskip
\noindent\textbf{Step 4: Bound \(\|V\|\) and \(\|VP_0\|\).}

Let \(Q_{\le T}^{(\mu)}\) and \(Q_{\le T}^{(\omega)}\) denote the orthogonal projectors onto the subspaces spanned by weights of depth at most \(T\) below the highest weight in \(\mathcal H_\mu\) and \(\mathcal H_\omega\), respectively.

Every vector in the total-weight slice \(W_{\nu-\delta}\) is a sum of tensors of the form
\[
(\mu-\delta_\mu)\otimes(\omega-\delta_\omega),
\qquad
\delta_\mu,\delta_\omega\in Q_+,
\qquad
\delta_\mu+\delta_\omega=\delta\ .
\]
In particular,
\[
|\delta_\mu|\le |\delta|,
\qquad
|\delta_\omega|\le |\delta|.
\]
Hence the slice projector factors through the shallow-depth subspaces:
\begin{equation}\label{eq:Q_depth_factorization}
    Q_{\nu-\delta}
    =
    Q_{\nu-\delta}\bigl(Q_{\le |\delta|}^{(\mu)}\otimes Q_{\le |\delta|}^{(\omega)}\bigr).
\end{equation}

Fix a positive root \(\alpha\). By the \(\alpha\)-string discussion in the preliminaries, the restrictions of \(\mathcal H_\mu\) and \(\mathcal H_\omega\) to the \(\mathfrak{sl}_2\)-subalgebra generated by \(E_{\pm\alpha}\) decompose into \(\alpha\)-strings. On the shallow subspaces \(Q_{\le |\delta|}^{(\mu)}\) and \(Q_{\le |\delta|}^{(\omega)}\), every vector lies at position \(q\le |\delta|\) in its local \(\alpha\)-string. Hence the operator norms of \(E_{\pm\alpha}\) on these restricted subspaces are controlled by \eqref{eq:alpha_string_norm_bound} once we bound the relevant length parameter \(m\).

For \(\mathcal H_\mu\), the length parameter of any \(\alpha\)-string is bounded by \eqref{eq:alpha_string_length_bound}:
\[
m\le \mu_1-\mu_d\le |\mu|=O(n).
\]
Therefore
\begin{equation}\label{eq:E_mu_bound}
    \bigl\|Q_{\le |\delta|}^{(\mu)}E_{\pm\alpha}^{(\mu)}Q_{\le |\delta|}^{(\mu)}\bigr\|
    =
    O\!\left(\sqrt{|\delta|\,n}\right).
\end{equation}

Similarly, for \(\mathcal H_\omega\),
\[
m\le \omega_1-\omega_d=\|\omega\|,
\]
and therefore
\begin{equation}\label{eq:E_omega_bound}
    \bigl\|Q_{\le |\delta|}^{(\omega)}E_{\pm\alpha}^{(\omega)}Q_{\le |\delta|}^{(\omega)}\bigr\|
    =
    O\!\left(\sqrt{|\delta| \|\omega\|}\right).
\end{equation}

Using \eqref{eq:Q_depth_factorization}, \eqref{eq:E_mu_bound}, and \eqref{eq:E_omega_bound}, and summing over the finitely many positive roots, we obtain
\begin{equation}\label{eq:V_norm_bound}
    \|V\|
    =
    O\!\left(|\delta|\,\sqrt{n\,\|\omega\|}\right).
\end{equation}

Since we use a tighter version of the Davis-Kahan, we want the tighter norm of \(\|VP_0\|\). Since \(\ket{\omega}\) is a highest-weight vector, it is annihilated by every raising operator:
\[
E_\alpha^{(\omega)}\ket{\omega}=0
\qquad
(\alpha>0).
\]
Hence only the terms \(E_\alpha^{(\mu)}\otimes E_{-\alpha}^{(\omega)}\) survive on \(P_0\), and
\begin{equation}
    VP_0 = 2Q_{\nu-\delta}\sum_{\alpha>0}
    \bigl(E_\alpha^{(\mu)}\otimes E_{-\alpha}^{(\omega)}\bigr)P_0\ .
\end{equation}
Moreover, on the highest-weight line of \(\mathcal H_\omega\), the vector \(\ket{\omega}\) sits at position \(q=0\) in every \(\alpha\)-string. Therefore, by \eqref{eq:alpha_string_minus},
\[
\|E_{-\alpha}^{(\omega)}\ket{\omega}\|^2
= \langle \omega,\alpha^\vee\rangle
\le \|\omega\|.
\]
Combining this with \eqref{eq:E_mu_bound} gives
\begin{equation}\label{eq:VP0_bound}
    \|VP_0\| =  O\!\left(\sqrt{n\,|\delta| \|\omega\|}\right).
\end{equation}


\medskip
\noindent\textbf{Step 5: Apply Davis--Kahan.}

We are ready to apply Lemma~\ref{lem:davis_kahan}:
$$ \|P-P_0\|\le \frac{\|VP_0\|}{\Delta-\|V\|} .$$
By \eqref{eq:V_norm_bound} and the assumption \(|\delta|^2\|\omega\|\ll n\),
\[
\|V\|=o(n).
\]
Together with \eqref{eq:Delta_def}, this implies
\[
\|V\|\ll \Delta.
\]
We can therefore apply Lemma~\ref{lem:davis_kahan} to the pair \(H_0\) and \(H=H_0+V\), with unperturbed projector \(P_0=\Pi_{\mu-\delta}^{(\omega)}\) and perturbed projector \(P=\widetilde{\Pi}_{\nu-\delta}\). Using~\eqref{eq:Delta_def} and \eqref{eq:VP0_bound}, we obtain
\begin{equation}
    \left\|\widetilde{\Pi}_{\nu-\delta}-\Pi_{\mu-\delta}^{(\omega)}\right\|\equiv\|P-P_0\|
    \le \frac{\|VP_0\|}{\Delta-\|V\|} 
    = \frac{O\!\left(\sqrt{n\,|\delta| \|\omega\|}\right)}{\Omega(n)} 
    = O\!\left(\sqrt{\frac{|\delta| \|\omega\|}{n}}\right).
\end{equation}
This proves the lemma.
\end{proof}

The perturbative estimate only controls shallow weight sectors at small $|\delta|$, so we need to bound the contribution of the lower-weight tail.

\begin{lem}[Tail mass]\label{lem:tail}
Let \(\rho\) be a density matrix of rank \(r\le d\) with strictly decreasing nonzero spectrum \(x_1>\cdots>x_r>0\).
Let \(\mu\) be a dominant highest weight of length at most~\(r\), and write the normalized irrep state as
\[
    \rho_\mu =
    \sum_{\delta\in Q_+} p_\mu(\delta)\,\Pi_{\mu-\delta},
    \qquad
    p_\mu(\delta)=\frac{x^{\mu-\delta}}{s_\mu(x)}\ .
\]
Define the tail mass
\begin{equation}
    \varepsilon_\mu(T) := 
    \sum_{|\delta|>T} p_\mu(\delta)\,m_\mu(\delta),
    \qquad
    m_\mu(\delta):= \dim W_\mu(\mu-\delta).
\end{equation}
If \(r=1\), then \(\varepsilon_\mu(T)=0\) for all \(T\).
If \(r\ge 2\), let
\[
    p:= \max_{1\le i\le r-1}\frac{x_{i+1}}{x_i}\in (0,1),
    \qquad
    \kappa:= \frac{r(r-1)}2-1.
\]
Then
\begin{equation}
    \varepsilon_\mu(T)=O\!\left(T^\kappa p^T\right).
\end{equation}
The implicit constant depends only on \(r\) and the spectrum \(x\).
\end{lem}
\begin{proof}
    See Appendix~\ref{app:proofs}.
\end{proof}


We now estimate the classical contribution (Bhattacharyya coefficient) to the fidelity~\eqref{eq:block_fidelity}. Since \(\nu=\mu+\omega\), the ratio 
\begin{equation}
    \frac{p_\nu(\delta)}{p_\mu(\delta)}=\frac{x^{\mu+\omega-\delta}/s_{\mu+\omega}(x)}{x^{\mu-\delta}/s_{\mu}(x)} = \frac{x^\omega s_\mu(x)}{s_{\mu+\omega}(x)}
\end{equation}
is \emph{independent} of \(\delta\), and it remains to show that this ratio is asymptotically close to \(1\).

\begin{lem}[Probability Ratio]\label{lem:ratio}
    Let \(\mu,\omega\) be two \(\mathrm{GL}(d,\mathbb C)\) irreps, and let \(x\) be a probability vector with \(r \le d\) strictly decreasing non-zero entries \(x_1>\cdots>x_r>0\). Let
    \[
    \xi := \min_{1 \le i \le r-1} (\ln x_i - \ln x_{i+1})
    \]
    be the logarithmic spectral gap. The ratio
    \begin{equation}
    Z(x):= \frac{x^\omega s_\mu(x)}{s_{\mu+\omega}(x)}
    \end{equation}
    satisfies
    \[
        1 \ge Z(x) \ge 1 - O\left(e^{-(g_\mu+1)\xi}\right)\ .
    \]
\end{lem}
\begin{proof}
    See Appendix~\ref{app:proofs}.
\end{proof}

We are now ready to complete the proof of the Theorem~\ref{thm:main}.

\begin{proof}[Proof of Theorem~\ref{thm:main} continued]
Picking up from~\eqref{eq:block_fidelity},
\begin{equation}
    F(\C_{\mu\to\nu}(\rho_\mu),\rho_\nu)
    \ge
    \sum_{\delta\in Q_+}
    \sqrt{p_\mu(\delta)p_\nu(\delta)}\,
    F(\mathcal K_\omega(\Pi_{\mu-\delta}),\widetilde{\Pi}_{\nu-\delta}) \ .
\end{equation}

For a fixed weight block \(\delta\), define
\begin{equation}
    F_{\mathrm{block}}(\delta)
    :=
    F(\mathcal K_\omega(\Pi_{\mu-\delta}),\widetilde{\Pi}_{\nu-\delta})
    =
    \Tr\sqrt{
        \widetilde{\Pi}_{\nu-\delta}
        \mathcal K_\omega(\Pi_{\mu-\delta})
        \widetilde{\Pi}_{\nu-\delta}
    } \ .
\end{equation}
Substituting~\eqref{eq:kraus_proj}, the operator under the square root becomes
\begin{equation}
    \frac{d_\mu}{d_\nu}\,
    \widetilde{\Pi}_{\nu-\delta}
    P_\nu
    \bigl(\Pi_{\mu-\delta}\otimes \ketbra{\omega}{\omega}\bigr)
    P_\nu
    \widetilde{\Pi}_{\nu-\delta}\ .
\end{equation}
Using~\eqref{eq:proj_mu_omega}, namely
\[
\Pi_{\mu-\delta}^{(\omega)} =
\Pi_{\mu-\delta}\otimes \ketbra{\omega}{\omega},
\]
the identity \(\widetilde{\Pi}_{\nu-\delta}=Q_{\nu-\delta}P_\nu=P_\nu Q_{\nu-\delta}\) and Lemma~\ref{lem:principle_angles}, the block fidelity simplifies to
\begin{equation}
    F_{\mathrm{block}}(\delta)
    =
    \sqrt{\frac{d_\mu}{d_\nu}}\,
    \Tr\sqrt{
        \widetilde{\Pi}_{\nu-\delta}
        \Pi_{\mu-\delta}^{(\omega)}
        \widetilde{\Pi}_{\nu-\delta}
    }
    =
    \sqrt{\frac{d_\mu}{d_\nu}}
    \sum_{i=1}^{m(\delta)} \cos\theta_i\ ,
\end{equation}
where \(\theta_i\) are the principal angles between
\(\mathrm{im}(\widetilde{\Pi}_{\nu-\delta})\) and
\(\mathrm{im}(\Pi_{\mu-\delta}^{(\omega)})\), and their rank is given by Lemma~\ref{lem:kostka}:
\[
m(\delta)=K_{\mu,\mu-\delta}=K_{\nu,\nu-\delta}\ .
\]

Using \(\cos\theta\ge 1-\sin^2\theta\), together with
\begin{equation}
    \sin\theta_{\max}=
    \left\|\widetilde{\Pi}_{\nu-\delta}-\Pi_{\mu-\delta}^{(\omega)}\right\|=
    O\!\left(\sqrt{\frac{|\delta|\|\omega\|}{n}}\right)
\end{equation}
from Lemma~\ref{lem:perturbation}, we obtain
\begin{equation}
    \sum_{i=1}^{m(\delta)} \cos\theta_i \ge
    m(\delta)\left(
        1-O\!\left(\frac{|\delta|\|\omega\|}{n}\right)
    \right).
\end{equation}
Therefore
\begin{equation}
    F_{\mathrm{block}}(\delta) \ge
    \sqrt{\frac{d_\mu}{d_\nu}}\,
    m(\delta)\left(
        1-O\!\left(\frac{|\delta|\|\omega\|}{n}\right)
    \right).
\end{equation}

We set a weight cutoff \( |\delta|\le T = O(n^\epsilon)\) for an arbitrarily small
\(\epsilon>0\), chosen so that \(T^2\|\omega\|\ll n\), as required in
Lemma~\ref{lem:perturbation}. Furthermore, we have $\|\omega\|\le|\omega|=d(\mu,\nu)$. Then the subspace-deviation error is uniformly bounded by
\(O(d(\mu,\nu)/n^{1-\epsilon})\). We have
\begin{equation}
    F_{\mathrm{block}}(\delta) \ge
    \sqrt{\frac{d_\mu}{d_\nu}}\,
    m(\delta)\left(
        1-O\left(\frac{d(\mu,\nu)}{n^{1-\epsilon}}\right)
    \right).
\end{equation}

For the dimension ratio $d_\mu/d_\nu$,  we have the following lemma.
\begin{lem}[Stability of the Weyl dimension]\label{lem:dim_ratio}
Let \(\mu,\nu\) be dominant highest weights of \(\mathrm{GL}(d)\) with
\(\nu=\mu+\omega\), where \(\omega\) is dominant. Assume that \(\mu\) has at most \(r\) nonzero rows, and that
\[
g_\mu=\min_{1\le i\le r-1}(\mu_i-\mu_{i+1})=\Theta(n),
\qquad
\mu_r=\Theta(n),
\qquad
|\omega|=o(n).
\]
Then
\begin{equation}
    \frac{d_\mu}{d_\nu}
    =
    1-O\!\left(\frac{|\omega|}{n}\right)
    =
    1-O\!\left(\frac{d(\mu,\nu)}{n}\right).
\end{equation}
\end{lem}
\begin{proof}
    See Appendix~\ref{app:proofs}.
\end{proof}


Substituting the perturbation estimate and the dimension ratio into~\eqref{eq:block_fidelity}, we obtain
\begin{equation}
    F(\C_{\mu\to\nu}(\rho_\mu),\rho_\nu)
    \ge
    \left(
        1-O\!\left(\frac{d(\mu,\nu)}{n^{1-\epsilon}}\right)
    \right)
    \sum_{|\delta|\le T}
    p_\mu(\delta)\,m(\delta)\,
    \sqrt{\frac{p_\nu(\delta)}{p_\mu(\delta)}} \ .
\end{equation}
By Lemma~\ref{lem:ratio}, the probability ratio satisfies
\[
\frac{p_\nu(\delta)}{p_\mu(\delta)}
=
1-O\!\left(e^{-\xi g_\mu}\right)
\]
uniformly in \(\delta\), for some \(\xi>0\) depending only on the spectral gap of \(\rho\).
By Lemma~\ref{lem:tail}, the omitted mass beyond depth \(T\) is
\(O(T^\kappa p^T)\), which is superpolynomially small when \(T=O(n^\epsilon)\). Since
\begin{equation}
    \sum_{\delta} p_\mu(\delta)m(\delta) 
    = \Tr\!\left(\sum_\delta p_\mu(\delta)\Pi_{\mu-\delta}\right)
    = \Tr \rho_\mu = 1\ ,
\end{equation}
we conclude
\begin{equation}
    F(\C_{\mu\to\nu}(\rho_\mu),\rho_\nu)
    \ge
    1-O\!\left(\frac{d(\mu,\nu)}{n^{1-\epsilon}}\right).
\end{equation}
This proves the forward-direction estimate.\\

Finally, we bound the \emph{reverse-direction} fidelity
\(F(\C_{\nu\to\mu}(\rho_\nu),\rho_\mu)\) using essentially the same strategy as for the forward direction.
Using Proposition~\ref{prop:reverse} together with the Stinespring form of \(\C_{\mu\to\nu}\), we obtain
\begin{equation}\label{eq:reverse_channel}
    \C_{\nu\to\mu}(X) =
    \frac{d_\nu}{d_\mu}\,\C_{\mu\to\nu}^{\dagger}(X) =
    \Tr_\omega\!\bigl(V X V^\dagger\bigr).
\end{equation}
To see this, we note that for every \(X\in \mathcal B(\mathcal H_\nu)\) and \(Y\in \mathcal B(\mathcal H_\mu)\),
\begin{multline}
    \Tr\!\left[X\,\C_{\mu\to\nu}(Y)\right]
    =\frac{d_\mu}{d_\nu}\Tr\!\left[X\,V^\dagger(Y\otimes \id_\omega)V\right] 
    \\
    =\frac{d_\mu}{d_\nu}\Tr\!\left[VXV^\dagger (Y\otimes \id_\omega)\right] 
    =\Tr\!\left[
        \frac{d_\mu}{d_\nu}\Tr_\omega(VXV^\dagger)\,Y
    \right],
\end{multline}
so \(\C_{\mu\to\nu}^{\dagger}(X)=\frac{d_\mu}{d_\nu}\Tr_\omega(VXV^\dagger)\), and hence~\eqref{eq:reverse_channel} follows.

Now define
\begin{equation}
    \rho_\mu^{(\omega)}
    :=
    \rho_\mu\otimes \ketbra{\omega}{\omega}
    =
    \sum_{\delta\in Q_+} p_\mu(\delta)\,\Pi_{\mu-\delta}^{(\omega)}.
\end{equation}
Since \(\Tr_\omega(\widetilde{\rho}_\nu)=\C_{\nu\to\mu}(\rho_\nu)\) where $\widetilde{\rho}_\nu=V\rho_\nu V^\dagger$ and
\(\Tr_\omega(\rho_\mu^{(\omega)})=\rho_\mu\), monotonicity of fidelity under the partial trace
\(\Tr_\omega\) gives
\begin{equation}\label{eq:reverse_reduce}
    F(\C_{\nu\to\mu}(\rho_\nu),\rho_\mu)
    =
    F\!\left(\Tr_\omega(\widetilde{\rho}_\nu),\Tr_\omega(\rho_\mu^{(\omega)})\right)
    \ge
    F(\widetilde{\rho}_\nu,\rho_\mu^{(\omega)}).
\end{equation}
Expanding both states in the same total-weight decomposition yields
\begin{equation}\label{eq:reverse_blocks}
    F(\widetilde{\rho}_\nu,\rho_\mu^{(\omega)})
    =
    \sum_{\delta\in Q_+}
    \sqrt{p_\mu(\delta)p_\nu(\delta)}\,
    F(\widetilde{\Pi}_{\nu-\delta},\Pi_{\mu-\delta}^{(\omega)}).
\end{equation}

For \(|\delta|\le T=O(n^\eps)\), exactly the same principal-angle estimate as above gives
\begin{equation}
    F(\widetilde{\Pi}_{\nu-\delta},\Pi_{\mu-\delta}^{(\omega)})
    =
    \sum_{i=1}^{m(\delta)} \cos\theta_i
    \ge
    m(\delta)\left(
        1-O\!\left(\frac{|\delta||\omega|}{n}\right)
    \right).
\end{equation}
Compared with the forward direction, the only difference is that the prefactor
\(\sqrt{d_\mu/d_\nu}\) is absent. Using the same cutoff \(T\), the same probability-ratio estimate from Lemma~\ref{lem:ratio}, and the same tail bound from Lemma~\ref{lem:tail}, we obtain
\begin{equation}
    F(\widetilde{\rho}_\nu,\rho_\mu^{(\omega)})
    \ge
    1-O\!\left(\frac{|\omega|}{n^{1-\epsilon}}\right).
\end{equation}
Combining this with~\eqref{eq:reverse_reduce} and recalling that
\(|\omega|=d(\mu,\nu)\), we conclude
\begin{equation}
    F(\C_{\nu\to\mu}(\rho_\nu),\rho_\mu)
    \ge
    1-O\!\left(\frac{d(\mu,\nu)}{n^{1-\epsilon}}\right).
\end{equation}
This proves the reverse-direction estimate.


\end{proof}


\section*{Acknowledgments}
    We would like to thank Daniel Bump, Jiani Fei and Cynthia Yan for discussions and insights about covariant maps.


\appendix

\section{Proofs}\label{app:proofs}
In this appendix, we provide the proofs for the propositions and lemmas used in the main text. 
\begin{mprop}{1}[Stinespring Representation]
    Let $V: \mathcal{H}_\nu \to \mathcal{H}_\mu \otimes \mathcal{H}_\omega$ be the isometric embedding of the irrep $\nu$ onto the tensor product of irrep $\mu$ and $\omega$, and these highest weights satisfy $\nu = \mu + \omega$. The covariant channel $\mathcal{C}_{\mu \to \nu}(\sigma_\mu) = \frac{d_\mu}{d_\nu} V^\dagger (\sigma_\mu\otimes \id_\omega) V$ has a Choi matrix $J$ that is proportional to the projector $\Pi_\omega$ onto the PRV component $(\nu - \mu)^+$ in $\nu \otimes \mu^*$. Specifically, $J = \frac{d_\mu}{d_\omega} \Pi_\omega$.
\end{mprop}

\begin{proof}
Let
\[
\ket{\Omega}=\sum_K \ket{\mu,K}\otimes\ket{\mu^*,K^*}
\]
be the unnormalized maximally entangled vector. By definition,
\[
J=(\mathcal C_{\mu\to\nu}\otimes \mathrm{id})(\ketbra{\Omega}{\Omega}),
\]
so using
\[
\mathcal C_{\mu\to\nu}(\sigma)
=\frac{d_\mu}{d_\nu}\,V^\dagger(\sigma\otimes I_\omega)V,
\]
we get
\[
J=\frac{d_\mu}{d_\nu}\sum_{K,L}\sum_W
\Bigl(
V^\dagger \ketbra{\mu,K;\omega,W}{\mu,L;\omega,W} V
\Bigr)
\otimes
\ketbra{\mu^*,K^*}{\mu^*,L^*}\ .
\]

Define a linear map
\[
\tilde V:\mathcal H_\nu\otimes \mathcal H_{\mu^*}\to \mathcal H_\omega
\]
by
\[
\tilde V\bigl(\ket{\nu,N}\otimes\ket{\mu^*,K^*}\bigr)
:=
(\bra{\mu,K}\otimes I_\omega)V\ket{\nu,N}\ .
\]
Then its matrix elements satisfy
\[
\bra{\nu,N;\mu^*,K^*}\tilde V^\dagger \tilde V\ket{\nu,N';\mu^*,L^*}
=
\sum_W
\bra{\nu,N}V^\dagger\ket{\mu,K;\omega,W}
\bra{\mu,L;\omega,W}V\ket{\nu,N'}.
\]
Comparing with the expression above for \(J\), we obtain
\[
J=\frac{d_\mu}{d_\nu}\,\tilde V^\dagger \tilde V.
\]

Now \(V\) is an intertwiner \(V:\nu\to \mu\otimes\omega\), so \(\tilde V\) is an intertwiner
\[
\tilde V:\nu\otimes \mu^* \to \omega.
\]
Hence \(\tilde V^\dagger\tilde V\) commutes with the representation \(\nu\otimes\mu^*\), and its support lies in the \(\omega\)-isotypic component.

At this point we use multiplicity one. Since \(\nu=\mu+\omega\), the irrep \(\nu\) is the Cartan component of \(\mu\otimes\omega\), so
\[
\dim \operatorname{Hom}_{U(d)}(\nu,\mu\otimes\omega)=1.
\]
Equivalently,
\[
\dim \operatorname{Hom}_{U(d)}(\nu\otimes\mu^*,\omega)=1.
\]
Therefore \(\omega\) occurs with multiplicity one in \(\nu\otimes\mu^*\). If \(\Pi_\omega\) denotes the orthogonal projector onto that unique copy, then
\[
\tilde V^\dagger\tilde V = c\,\Pi_\omega
\]
for some \(c>0\).

To determine \(c\), take the trace:
\[
\operatorname{Tr}(\tilde V^\dagger\tilde V)
=
\sum_{N,K}\left\|
(\bra{\mu,K}\otimes I_\omega)V\ket{\nu,N}
\right\|^2
=
\sum_N \|V\ket{\nu,N}\|^2
=
\operatorname{Tr}(V^\dagger V)
=
d_\nu,
\]
since \(V\) is an isometry. On the other hand,
\[
\operatorname{Tr}(\tilde V^\dagger\tilde V)=c\,\operatorname{Tr}(\Pi_\omega)=c\,d_\omega.
\]
Thus
\[
c=\frac{d_\nu}{d_\omega},
\qquad\text{so}\qquad
\tilde V^\dagger\tilde V=\frac{d_\nu}{d_\omega}\Pi_\omega.
\]
Substituting back into the Choi formula gives
\[
J=\frac{d_\mu}{d_\nu}\tilde V^\dagger\tilde V
=\frac{d_\mu}{d_\omega}\Pi_\omega.
\]

Finally, since under the assumption \(\nu=\mu+\omega\) we have
\[
(\nu-\mu)^+=\omega,
\]
this is exactly the projector onto the PRV component of \(\nu\otimes\mu^*\).
\end{proof}


\begin{mprop}{2}[Reverse Cloner]
Let $\C_{\mu\to\nu}:\mathcal B(\mathcal H_\mu)\to \mathcal B(\mathcal H_\nu)$
be the generalizer cloning map, and $J_{\mu\to\nu}=\frac{d_\mu}{d_\omega}\Pi_\omega$ be its Choi operator. Let $\sigma_\mu :=  \id_\mu/d_\mu$.
Then the Petz recovery map of \(\C_{\mu\to\nu}\) with respect to \(\sigma_\mu\) is exactly the reverse  cloning map:
\[
\mathcal R_{\sigma_\mu,\C_{\mu\to\nu}}=\C_{\nu\to\mu}\ .
\]
Equivalently,
\[
\C_{\nu\to\mu}(X)=\frac{d_\nu}{d_\mu}\,\C_{\mu\to\nu}^{\dagger}(X),
\]
where \(\C_{\mu\to\nu}^{\dagger}\) is the adjoint map.
\end{mprop}

\begin{proof}
We first observe that \(\C_{\mu\to\nu}\) sends the maximally mixed state to the maximally mixed state. Indeed, since \(\C_{\mu\to\nu}\) is \(\mathrm U(d)\)-covariant,
\[
\C_{\mu\to\nu}\!\left(U_\mu(g)\sigma_\mu U_\mu(g)^\dagger\right)
=
U_\nu(g)\C_{\mu\to\nu}(\sigma_\mu)U_\nu(g)^\dagger
\qquad \forall g\in \mathrm U(d).
\]
Because \(\sigma_\mu=\id_\mu/d_\mu\) is invariant and \(U_\nu\) is irreducible, Schur's lemma implies that \(\C_{\mu\to\nu}(\sigma_\mu)\) is proportional to the identity on \(\mathcal H_\nu\). Since it has unit trace,
\[
\C_{\mu\to\nu}(\sigma_\mu)=\sigma_\nu:= \frac{\id_\nu}{d_\nu}.
\]

Therefore the Petz map with respect to \(\sigma_\mu\) is
\[
\mathcal R_{\sigma_\mu,\C_{\mu\to\nu}}(X)
=
\sigma_\mu^{1/2}\,
\C_{\mu\to\nu}^{\dagger}
\!\left(
\sigma_\nu^{-1/2}X\sigma_\nu^{-1/2}
\right)
\,\sigma_\mu^{1/2}
=
\frac{d_\nu}{d_\mu}\,\C_{\mu\to\nu}^{\dagger}(X).
\]
So it remains to identify the scaled adjoint with the reverse cloning channel.

Let
\[
\operatorname{vec}:\operatorname{Hom}(\mathcal H_\mu,\mathcal H_\nu)\longrightarrow \mathcal H_\mu^*\otimes \mathcal H_\nu
\]
be the canonical vectorization isometry. Choose an orthonormal basis
\[
\{\psi_a\}_{a=1}^{d_\omega}
\subset \mathcal H_\mu^*\otimes \mathcal H_\nu
\]
for the range of \(\Pi_\omega\), and define Kraus operators \(K_a\in \operatorname{Hom}(\mathcal H_\mu,\mathcal H_\nu)\) by
\[
\psi_a=\operatorname{vec}(K_a).
\]
Then
\[
J_{\mu\to\nu}
=
\frac{d_\mu}{d_\omega}
\sum_{a=1}^{d_\omega}
\ketbra{\psi_a}{\psi_a},
\]
and hence
\[
\C_{\mu\to\nu}(X)
=
\frac{d_\mu}{d_\omega}
\sum_{a=1}^{d_\omega}
K_a X K_a^\dagger.
\]
Consequently,
\[
\frac{d_\nu}{d_\mu}\,\C_{\mu\to\nu}^{\dagger}(Y)
=
\frac{d_\nu}{d_\omega}
\sum_{a=1}^{d_\omega}
K_a^\dagger Y K_a.
\]

Now the span
\[
\mathsf K_\omega:= \operatorname{span}\{K_a\}
\subset \operatorname{Hom}(\mathcal H_\mu,\mathcal H_\nu)
\]
carries the irrep \(\omega\) under the action
\[
K\longmapsto U_\nu(g) K U_\mu(g)^\dagger.
\]
Taking Hilbert--Schmidt adjoints defines an antiunitary intertwiner
\[
K\longmapsto K^\dagger
\]
from \(\operatorname{Hom}(\mathcal H_\mu,\mathcal H_\nu)\) to \(\operatorname{Hom}(\mathcal H_\nu,\mathcal H_\mu)\). Therefore
\[
\mathsf K_{\omega^*}:= \operatorname{span}\{K_a^\dagger\}
\subset \operatorname{Hom}(\mathcal H_\nu,\mathcal H_\mu)
\]
carries the dual irrep \(\omega^*\). Since
\[
\omega^*=(-w_0\omega)=(\mu-\nu)^+,
\]
this is exactly the PRV component for the reversed pair \((\nu,\mu)\).

Because \(\{K_a^\dagger\}_{a=1}^{d_\omega}\) is again an orthonormal family in Hilbert--Schmidt norm, the Choi operator of the map
\[
Y\longmapsto \frac{d_\nu}{d_\omega}\sum_{a=1}^{d_\omega}K_a^\dagger Y K_a
\]
is precisely
\[
\frac{d_\nu}{d_{\omega^*}}\Pi_{\omega^*},
\]
where \(\Pi_{\omega^*}\) denotes the projector onto the reverse PRV component
\[
\omega^*=(\mu-\nu)^+\subset \nu^*\otimes \mu.
\]
Since \(d_{\omega^*}=d_\omega\), this is exactly the cloning Choi operator for the reversed channel \(\C_{\nu\to\mu}\). Hence
\[
\frac{d_\nu}{d_\mu}\,\C_{\mu\to\nu}^{\dagger}
=
\C_{\nu\to\mu}.
\]
Combining this with the Petz formula above proves the claim.
\end{proof}


\begin{comment}
\begin{mprop}{3}[Optimality of the Generalized Cloning Map]
Let $\mu$ and $\nu$ be irreps of GL$(d,\mathbb C)$, and $\rho_\mu$ and $\rho_\mu$ be corresponding normalized irreps of $\rho$. Among all covariant quantum channels $\mathcal{N}: \mathcal{B}(\mathcal{H}_\mu) \to \mathcal{B}(\mathcal{H}_\nu)$, the generalized cloning map $\mathcal{C}_{\mu \to \nu}$ defined by its Choi operator $J_\omega = \frac{d_\mu}{d_\omega} \Pi_\omega$, where $\omega := (\nu - \mu)^+$, maximizes the overlap $\Tr\,(\mathcal{N}(\rho_\mu)\rho_\nu)$.
\end{mprop}

\begin{proof}
    By the Choi-Jamio\l{}kowski isomorphism and Schur's Lemma, any $\mathrm{U}(d)$-covariant channel $\mathcal{N}$ corresponds to a Choi operator $J_\mathcal{N}$ that decomposes into a convex combination of orthogonal projectors onto the irreps $\lambda$ in the Littlewood-Richardson decomposition of $\mu^* \otimes \nu$:
\begin{equation}
    J_\mathcal{N} = \sum_{\lambda \in \mu^* \otimes \nu} c_\lambda \frac{d_\mu}{d_\lambda} \Pi_\lambda, \quad \sum_\lambda c_\lambda = 1, \quad c_\lambda \ge 0\ .
\end{equation}
The trace overlap between the channel output and the target state is:
\begin{equation}
    O(\mathcal{N}) = \Tr\big[ \mathcal{N}(\rho_\mu) \rho_\nu \big] = \Tr\big[ J_\mathcal{N} (\rho_\mu^\top \otimes \rho_\nu) \big]\ .
\end{equation}
Let $X = \mathrm{diag}(x_1, \dots, x_d)$ be the diagonal matrix of eigenvales of $\rho$. The irrep states are given by $\rho_\mu = \frac{\pi_\mu(X)}{s_\mu(x)}$ and $\rho_\nu = \frac{\pi_\nu(X)}{s_\nu(x)}$, where $s_\lambda(x)$ is the Schur polynomial. Because transposition maps the representation to its dual, $\rho_\mu^\top \propto \pi_{\mu^*}(X)$. The tensor product is exactly the joint representation evaluated on the spectrum $X$:
\begin{equation}
    \rho_\mu^\top \otimes \rho_\nu = \frac{1}{s_\mu(x)s_\nu(x)} \pi_{\mu^* \otimes \nu}(X)\ .
\end{equation}
Substituting the convex decomposition of $J_\mathcal{N}$ yields:
\begin{equation}
    O(\mathcal{N}) = \frac{d_\mu}{s_\mu(x)s_\nu(x)} \sum_\lambda c_\lambda \frac{1}{d_\lambda} \Tr\big[ \Pi_\lambda \, \pi_{\mu^* \otimes \nu}(X) \big]\ .
\end{equation}
Because $\Pi_\lambda$ projects exactly onto the invariant subspace $\mathcal{H}_\lambda$, the trace over this subspace evaluates exactly to the character of $\lambda$, which is the Schur polynomial $s_\lambda(x)$:
\begin{equation}
    O(\mathcal{N}) = \frac{d_\mu}{s_\mu(x)s_\nu(x)} \sum_\lambda c_\lambda \left( \frac{s_\lambda(x)}{d_\lambda} \right)\ .
\end{equation}
To maximize this convex combination, we must place unit weight on the representation $\lambda$ that maximizes the normalized character $r_\lambda(x) := \frac{s_\lambda(x)}{d_\lambda}$. 
By the properties of highest-weight representations, every irrep $\lambda$ in the decomposition must physically contain the weight $\nu - \mu$. Therefore, its highest weight must dominate the PRV component $\omega = (\nu - \mu)^+$ in the standard dominance ordering: $\lambda \succeq \omega$.
For any decreasing sequence $x_1 \ge x_2 \ge \dots \ge x_d \ge 0$, the normalized Schur polynomial $r_\lambda(x)$ is monotonically decreasing with respect to the dominance order~\cite{cuttler2011inequalities,sra2016inequalities}. Since $\lambda \succeq \omega$ for all constituents, we have $\frac{s_\lambda(x)}{d_\lambda} \le \frac{s_\omega(x)}{d_\omega}$. Thus, the trace overlap is globally maximized by setting $c_\omega = 1$, corresponding uniquely to the map $\mathcal{C}_{\mu \to \nu}$ of our choice.
\end{proof}

    
\end{comment}




\begin{mprop}{3}[Commutativity] 
Let $\rho \in \mathcal{D}(\mathbb{C}^d)$ be a density matrix. Let $\mu$ and $\nu$ be irreps of $\mathrm{GL}(d,\mathbb C)$, $ \rho_\mu$ and $ \rho_\nu$ be the normalized $\mathrm{GL}(d,\mathbb C)$ irreps of $\rho$.
Let $\C_{\mu \to \nu}: \mathcal{B}(\h_\mu) \to \mathcal{B}(\h_\nu)$ be a quantum channel that is $\mathrm{U}(d)$-covariant, i.e., for all $g \in \mathrm{U}(d)$:
\begin{equation}
    \C_{\mu \to \nu}( U_\mu(g) X U_\mu(g)^{-1} ) = U_\nu(g) \C_{\mu \to \nu}(X) U_\nu(g)^{-1}
\end{equation}
where $U_\mu(g) := \pi_\mu(g)$ for $g \in \mathrm{U}(d)$.
Then, we have
\begin{equation}
    \left[ \C_{\mu \to \nu}(\rho_\mu), \rho_\nu \right] = 0 \ .
\end{equation}
\end{mprop}

\begin{proof}
The commutation relation is basis-independent, so we may work in the eigenbasis of the seed state $\rho$. Let $V \in \mathrm{U}(d)$ be the unitary that diagonalizes $\rho$ such that $\rho = \Lambda$. Consequently, the representation states $\rho_\mu$ and $\rho_\nu$ are images of the diagonal matrix $\Lambda$:
\begin{equation}
    \rho_\mu \propto \pi_\mu(\Lambda), \quad \rho_\nu \propto \pi_\nu(\Lambda) \ .
\end{equation}

Let $T \subset \mathrm{U}(d)$ be the maximal torus consisting of diagonal unitary matrices. Since $\Lambda$ is diagonal, it commutes with all elements of the torus: $[\rho_\mu, U_\mu(t)] = 0$ for all $t \in T$.

Applying the $\mathrm{U}(d)$-covariance property to an element $t \in T$:
\begin{equation}
    \C_{\mu\to\nu}( U_\mu(t) \rho_\mu U_\mu(t)^{-1} ) = U_\nu(t) \C_{\mu\to\nu}(\rho_\mu) U_\nu(t)^{-1} \ .
\end{equation}
Because $\rho_\mu$ commutes with $U_\mu(t)$, the left-hand side reduces directly to $\C_{\mu\to\nu}(\rho_\mu)$, which implies:
\begin{equation}
    [\C_{\mu\to\nu}(\rho_\mu), U_\nu(t)] = 0 \quad \forall t \in T \ .
\end{equation}

The representation space $\h_\nu$ decomposes orthogonally into weight spaces $\h_\nu = \bigoplus_{\delta} W_\nu(\delta)$. The $\mathrm{U}(d)$ torus action distinguishes these orthogonal subspaces via unique phase characters. Because $\C_{\mu\to\nu}(\rho_\mu)$ commutes with the entire torus $T$, Schur's Lemma dictates that it must be block-diagonal with respect to these weight spaces.

Simultaneously, the target state $\rho_\nu \propto \pi_\nu(\Lambda)$ acts purely as a scalar multiple of the identity on each distinct weight space $W_\nu(\delta)$. Because $\C_{\mu\to\nu}(\rho_\mu)$ is block-diagonal across the weight spaces and $\rho_\nu$ is a block-scalar on those identical spaces, the two operators must commute:
\begin{equation}
    [\C_{\mu\to\nu}(\rho_\mu), \rho_\nu] = 0 \ .
\end{equation}
\end{proof}




\begin{comment}
\begin{mprop}{4}[Variational Characterization]
Let $\mathcal N:\mathcal B(\mathcal H_\mu)\to \mathcal B(\mathcal H_\nu)$
be a \(\mathrm U(d)\)-covariant quantum channel, and let \(J_\mathcal N\) be its Choi operator. Let
$\omega:=(\nu-\mu)^+$ be the PRV component in \(\mu^*\otimes \nu\), and let \(C_2\) denote the quadratic Casimir operator of the representation \(U_\mu^*\otimes U_\nu\) acting on \(\mathcal H_\mu^*\otimes \mathcal H_\nu\). Define the average Casimir of the channel as
\[
K(\mathcal N):= \frac{1}{d_\mu}\Tr(C_2 J_\mathcal N),
\qquad
d_\mu:= \dim \mathcal H_\mu.
\]
Then
\[
K(\mathcal N)\ge C_2(\omega),
\]
with equality if and only if $N=\C_{\mu\to\nu}$.
Equivalently, the PRV channel is the unique \(\mathrm U(d)\)-covariant channel of minimal average Casimir.
\end{mprop}


\begin{proof}
Write the Schur-Weyl decomposition
\[
\mathcal H_\mu^*\otimes \mathcal H_\nu
\cong
\bigoplus_{\lambda} M_\lambda\otimes V_\lambda,
\]
where \(V_\lambda\) is the \(\lambda\)-irrep and \(M_\lambda\) is its multiplicity space. By covariance, the Choi operator of \(\mathcal N\) has the block form
\[
J_\mathcal N=\bigoplus_\lambda A_\lambda\otimes I_{V_\lambda},
\qquad
A_\lambda\ge 0.
\]
Since \(\mathcal N\) is trace-preserving,
\[
\Tr J_\mathcal N = d_\mu.
\]
Define
\[
p_\lambda:= \frac{d_\lambda\,\Tr A_\lambda}{d_\mu},
\qquad
d_\lambda:= \dim V_\lambda.
\]
Then \(p_\lambda\ge 0\) and
\[
\sum_\lambda p_\lambda=1.
\]

Because the quadratic Casimir acts by the scalar \(C_2(\lambda)\) on the irrep \(V_\lambda\), we obtain
\[
K(\mathcal N)
=
\frac{1}{d_\mu}\Tr(C_2 J_\mathcal N)
=
\sum_\lambda p_\lambda\,C_2(\lambda).
\]

Now \(\omega\) is the unique dominance-minimal highest weight appearing in \(\mu^*\otimes \nu\). Hence every other summand is of the form
\[
\lambda=\omega+\beta,
\qquad
\beta\in Q_+,
\qquad
\beta\neq 0.
\]
Using the highest-weight formula for the Casimir eigenvalue,
\[
C_2(\lambda)=(\lambda,\lambda+2\rho),
\]
we get
\[
C_2(\lambda)-C_2(\omega)
=
(\omega+\beta,\omega+\beta+2\rho)-(\omega,\omega+2\rho)
=
2(\omega+\rho,\beta)+(\beta,\beta)>0.
\]
Therefore
\[
C_2(\lambda)>C_2(\omega)
\qquad
\text{for every }\lambda\neq \omega
\text{ appearing in }\mu^*\otimes \nu.
\]
It follows that
\[
K(\mathcal N)=\sum_\lambda p_\lambda C_2(\lambda)\ge C_2(\omega),
\]
with equality if and only if
\[
p_\lambda=0
\qquad
\forall \lambda\neq \omega.
\]
So equality holds if and only if \(J_\mathcal N\) is supported entirely on the PRV block.

Since the PRV component occurs with multiplicity one, the \(\omega\)-block is one-dimensional on the multiplicity side, so on that block
\[
J_\mathcal N = a\,\Pi_\omega
\]
for some \(a>0\). The trace constraint \(\Tr J_\mathcal N=d_\mu\) then forces
\[
a\,d_\omega=d_\mu,
\qquad\text{hence}\qquad
J_\mathcal N=\frac{d_\mu}{d_\omega}\Pi_\omega.
\]
This is exactly the Choi operator of the PRV channel \(\C_{\mu\to\nu}\). Hence equality holds if and only if
\[
\mathcal N=\C_{\mu\to\nu}.
\]
\end{proof}
\end{comment}

\begin{mlemma}{3}[Kostka Numbers]
    Let $\mu,\omega$ be dominant integral highest weights of $\mathrm{GL}(d,\mathbb C)$, let $\delta\in Q_+$ be a weight offset, and let $g_\mu:= \min_{1\le i\le d-1}(\mu_i-\mu_{i+1})$ be the edge gap of $\mu$. Then the inner degeneracies (Kostka numbers) satisfy
    \begin{equation}
        K_{\mu,\mu-\delta} \le K_{\mu+\omega,\mu+\omega-\delta}.
    \end{equation}
    Moreover, equality holds when $|\delta|\le g_\mu$.
\end{mlemma}

\begin{proof}

The weight multiplicity of the $\mathrm{GL}(d,\mathbb C)$ irrep $\h_\lambda$ at weight $\beta$ is counted by Gelfand--Tsetlin (GT) patterns with top row $\lambda$ and weight $\beta$. More precisely, if $\mathrm{GT}(\lambda,\beta)$ denotes the set of triangular arrays
\begin{equation}
\begin{matrix}
m_{d,1} && m_{d,2} && \cdots && m_{d,d}\\
& m_{d-1,1} && \cdots && m_{d-1,d-1} & \\
&& \ddots && \ddots && \\
&&& m_{1,1} &&&
\end{matrix}
\end{equation}
satisfying
\begin{equation}
    m_{i+1,j}\ge m_{i,j}\ge m_{i+1,j+1}
    \qquad (1\le j\le i<d),
\end{equation}
with top row
\begin{equation}
    (m_{d,1},\dots,m_{d,d})=\lambda,
\end{equation}
and weight determined by
\begin{equation}
    \beta_i =
    \sum_{j=1}^i m_{i,j} -
    \sum_{j=1}^{i-1} m_{i-1,j},
    \qquad 1\le i\le d,
\end{equation}
then
\begin{equation}
    |\mathrm{GT}(\lambda,\beta)|=\dim W_\lambda(\beta)=K_{\lambda,\beta}.
\end{equation}

Take any pattern $M=(m_{i,j})\in \mathrm{GT}(\mu,\mu-\delta)$.
Define a new array $M'=(m'_{i,j})$ by
\begin{equation}
    m'_{i,j}:= m_{i,j}+\omega_j,
    \qquad 1\le j\le i\le d.
\end{equation}
We claim that $M'\in \mathrm{GT}(\mu+\omega,\mu+\omega-\delta)$.

First, its top row is
\begin{equation}
    m'_{d,j}=m_{d,j}+\omega_j=\mu_j+\omega_j,
\end{equation}
so the top row is indeed $\mu+\omega$.

Second, the interlacing conditions are preserved. The left inequality is immediate:
\begin{equation}
    m'_{i+1,j}-m'_{i,j}
    =
    (m_{i+1,j}+\omega_j)-(m_{i,j}+\omega_j)
    =
    m_{i+1,j}-m_{i,j}
    \ge 0.
\end{equation}
For the right inequality, we use that $\omega$ is dominant, so $\omega_j\ge \omega_{j+1}$:
\begin{equation}
    m'_{i,j}-m'_{i+1,j+1}
    =
    (m_{i,j}+\omega_j)-(m_{i+1,j+1}+\omega_{j+1})
    =
    (m_{i,j}-m_{i+1,j+1})+(\omega_j-\omega_{j+1})
    \ge 0.
\end{equation}
Hence $M'$ is again a valid Gelfand--Tsetlin pattern.

Third, its weight is shifted by $\omega$. If $\beta$ is the weight of $M$, then
\begin{equation}
\begin{aligned}
    \beta_i'
    &=
    \sum_{j=1}^i m'_{i,j}
    -
    \sum_{j=1}^{i-1} m'_{i-1,j} \\
    &=
    \sum_{j=1}^i (m_{i,j}+\omega_j)
    -
    \sum_{j=1}^{i-1} (m_{i-1,j}+\omega_j) \\
    &=
    \left(\sum_{j=1}^i m_{i,j}-\sum_{j=1}^{i-1} m_{i-1,j}\right)+\omega_i
    = \beta_i+\omega_i.
\end{aligned}
\end{equation}
Since $\beta=\mu-\delta$, this gives
\begin{equation}
    \beta'=(\mu-\delta)+\omega=\mu+\omega-\delta.
\end{equation}

Thus we have defined an injective map
\begin{equation}
    \mathrm{GT}(\mu,\mu-\delta)\hookrightarrow
    \mathrm{GT}(\mu+\omega,\mu+\omega-\delta),
    \qquad
    M\mapsto M'.
\end{equation}
Therefore
\begin{equation}
    K_{\mu,\mu-\delta}\le K_{\mu+\omega,\mu+\omega-\delta}.
\end{equation}\\


Now we prove the equality when $|\delta|\le g_\mu$. Let $M(\lambda)$ be the Verma module of highest weight $\lambda$.
Its weight space at $\lambda-\delta$ has dimension equal to Kostant's partition function, which we denote by $K(\delta)$:
\begin{equation}
    \dim M(\lambda)_{\lambda-\delta}=K(\delta).
\end{equation}
In particular, this dimension is independent of $\lambda$.

The irreducible module $\h_\lambda$ is the quotient of $M(\lambda)$ by the submodule generated by the simple singular vectors
\begin{equation}
    E_{-\alpha_i}^{\langle \lambda,\alpha_i^\vee\rangle+1}v_\lambda,
    \qquad i=1,\dots,d-1.
\end{equation}
These singular vectors first appear at depths
\begin{equation}
    \langle \lambda,\alpha_i^\vee\rangle+1
    =
    (\lambda_i-\lambda_{i+1})+1.
\end{equation}
Hence every nonzero vector in the kernel of the quotient map
\begin{equation}
M(\lambda)\twoheadrightarrow \h_\lambda
\end{equation}
has height at least
\begin{equation}
    g_\lambda+1,
    \qquad
    g_\lambda:= \min_i(\lambda_i-\lambda_{i+1}).
\end{equation}
Therefore, for every offset $\delta$ with $|\delta|\le g_\lambda$, the quotient map is an isomorphism on the weight space at $\lambda-\delta$, and
\begin{equation}
    K_{\lambda,\lambda-\delta}=K(\delta).
\end{equation}

Apply this first with $\lambda=\mu$. If $|\delta|\le g_\mu$, then
\begin{equation}
    K_{\mu,\mu-\delta}=K(\delta).
\end{equation}
Next apply it with $\lambda=\mu+\omega$. Since $\omega$ is dominant,
\begin{equation}
    (\mu_i+\omega_i)-(\mu_{i+1}+\omega_{i+1})
    \ge
    \mu_i-\mu_{i+1},
\end{equation}
so
\begin{equation}
    g_{\mu+\omega}\ge g_\mu.
\end{equation}
Hence the same assumption $|\delta|\le g_\mu$ implies $|\delta|\le g_{\mu+\omega}$, and therefore
\begin{equation}
    K_{\mu+\omega,\mu+\omega-\delta}=K(\delta).
\end{equation}
Combining the two identities yields
\begin{equation}
    K_{\mu,\mu-\delta}=K_{\mu+\omega,\mu+\omega-\delta},
\end{equation}
which proves the lemma.
\end{proof}


\begin{mlemma}{4}[Principle Angles]
Let $P$ and $Q$ be orthogonal projectors, and let $\theta_1,\dots,\theta_r$ denote the principal angles between $\operatorname{Im}(P)$ and $\operatorname{Im}(Q)$, where $r=\min(\operatorname{rank}P,\operatorname{rank}Q)$. Then
\begin{equation}
    \Tr \sqrt{PQP} = \sum_{i=1}^r \cos\theta_i \ .
\end{equation}
\end{mlemma}

\begin{proof}
Let $U$ and $V$ be matrices whose columns form orthonormal bases for $\operatorname{Im}(P)$ and $\operatorname{Im}(Q)$, respectively. Because $P$ and $Q$ are orthogonal projectors, they can be expressed as $P = UU^\dagger$ and $Q = VV^\dagger$.

By definition, the principal angles $\theta_1, \dots, \theta_r$ between these subspaces are determined by the singular values of the cross-projector matrix $U^\dagger V$. Specifically, the nonzero singular values $\sigma_i$ of $U^\dagger V$ satisfy $\sigma_i = \cos\theta_i$ for $i = 1, \dots, r$.

To determine the eigenvalues of $PQP$, we substitute the basis representations:
\begin{equation}
    PQP = (UU^\dagger)(VV^\dagger)(UU^\dagger) = U(U^\dagger V)(V^\dagger U)U^\dagger = U(U^\dagger V)(U^\dagger V)^\dagger U^\dagger \ .
\end{equation}
Let $M = U^\dagger V$. The operator $PQP$ takes the form $UMM^\dagger U^\dagger$. Because the columns of $U$ are orthonormal (i.e., $U^\dagger U = I$), the nonzero eigenvalues of $UMM^\dagger U^\dagger$ are identical to the nonzero eigenvalues of $MM^\dagger$. 

The eigenvalues of $MM^\dagger$ are simply the squared singular values of $M$. Therefore, the nonzero eigenvalues of $PQP$ are exactly
\begin{equation}
    \sigma_1^2, \dots, \sigma_r^2 = \cos^2\theta_1, \dots, \cos^2\theta_r \ .
\end{equation}
Because $PQP$ is positive semidefinite, the eigenvalues of its unique positive semidefinite square root, $\sqrt{PQP}$, are the principal square roots of these values:
\begin{equation}
    \cos\theta_1, \dots, \cos\theta_r \ ,
\end{equation}
together with possibly additional zeros if $\operatorname{rank}P > r$. Taking the trace, which is the sum of all eigenvalues, gives
\begin{equation}
    \Tr \sqrt{PQP} = \sum_{i=1}^r \cos\theta_i \ ,
\end{equation}
as claimed.
\end{proof}


\begin{mlemma}{5}[Monotonicity]
Let $\mathcal{N}$ be a quantum channel and $\rho = \mathcal{N}(\eta)$ for some input state $\eta \geq 0$. Let $K$ be a Kraus operator of $\mathcal{N}$ such that $\mathcal{N}(\cdot) = K(\cdot)K^\dagger + \sum_{i} L_i (\cdot) L_i^\dagger$, and define the unnormalized state $\rho' = K \eta K^\dagger$. For any density operator $\sigma$, the fidelity satisfies:
\begin{equation}
    F(\rho', \sigma) \leq F(\rho, \sigma) \ .
\end{equation}
\end{mlemma}

\begin{proof}
From the definition of the quantum channel, we can express the output state as $\rho = \rho' + \rho_{\text{rem}}$, where $\rho_{\text{rem}} = \sum_i L_i \eta L_i^\dagger$. Because the input state is positive semi-definite ($\eta \geq 0$) and the map $\sum_i L_i (\cdot) L_i^\dagger$ is completely positive, it follows that $\rho_{\text{rem}} \geq 0$. This gives the operator inequality:
\begin{equation}
    \rho \geq \rho' \ .
\end{equation}
Conjugating both sides of this inequality by the positive semi-definite operator $\sqrt{\sigma}$ preserves the operator ordering:
\begin{equation}
    \sqrt{\sigma} \rho \sqrt{\sigma} \geq \sqrt{\sigma} \rho' \sqrt{\sigma} \ .
\end{equation}
By the L\"owner-Heinz theorem, the square root function $f(x) = \sqrt{x}$ is operator monotone. Applying this to our inequality yields:
\begin{equation}
    \sqrt{\sqrt{\sigma} \rho \sqrt{\sigma}} \geq \sqrt{\sqrt{\sigma} \rho' \sqrt{\sigma}} \ .
\end{equation}
The trace operation is a positive linear functional, meaning it preserves operator inequalities. Taking the trace of both sides, we recover the definition of the quantum fidelity:
\begin{equation}
    \text{Tr}\left(\sqrt{\sqrt{\sigma} \rho \sqrt{\sigma}}\right) \geq \text{Tr}\left(\sqrt{\sqrt{\sigma} \rho' \sqrt{\sigma}}\right) \ ,
\end{equation}
which simplifies directly to:
\begin{equation}
    F(\rho, \sigma) \geq F(\rho', \sigma) \ .
\end{equation}
This concludes the proof.
\end{proof}


\begin{mlemma}{8}[Tail mass]
Let \(\rho\) be a density matrix of rank \(r\le d\) with strictly decreasing nonzero spectrum \(x_1>\cdots>x_r>0\).
Let \(\mu\) be a dominant highest weight of length at most \(r\), and write the normalized irrep state as
\[
    \rho_\mu
    =
    \sum_{\delta\in Q_+} p_\mu(\delta)\,\Pi_{\mu-\delta},
    \qquad
    p_\mu(\delta)=\frac{x^{\mu-\delta}}{s_\mu(x)},
\]
Define the tail mass
\begin{equation}
    \varepsilon_\mu(T)
    := 
    \sum_{|\delta|>T} p_\mu(\delta)\,m_\mu(\delta),
    \qquad
    m_\mu(\delta):= \dim W_\mu(\mu-\delta).
\end{equation}
If \(r=1\), then \(\varepsilon_\mu(T)=0\) for all \(T\).
If \(r\ge 2\), let
\begin{equation}
    p:= \max_{1\le i\le r-1}\frac{x_{i+1}}{x_i}\in (0,1),
    \qquad
    \kappa:= \frac{r(r-1)}2-1.
\end{equation}
Then
\begin{equation}
    \varepsilon_\mu(T)=O\!\left(T^\kappa p^T\right).
\end{equation}
The implicit constant depends only on \(r\) and the spectrum \(x\).
\end{mlemma}

\begin{proof}
The case \(r=1\) is trivial: there are no positive roots, hence no nonzero offsets \(\delta\), so \(\rho_\mu\) is supported entirely on the highest-weight space and \(\varepsilon_\mu(T)=0\).

Assume from now on that \(r\ge 2\).
Write the offset in the simple-root basis as
\begin{equation}
    \delta=\sum_{i=1}^{r-1} c_i \alpha_i,
    \qquad c_i\in \mathbb Z_{\ge 0},
    \qquad |\delta|=\sum_{i=1}^{r-1} c_i.
\end{equation}
Since \(\alpha_i=e_i-e_{i+1}\), we have
\begin{equation}
    x^{\mu-\delta}
    =
    x^\mu \prod_{i=1}^{r-1}\left(\frac{x_{i+1}}{x_i}\right)^{c_i}
    \le
    x^\mu\, p^{|\delta|}.
\end{equation}
Moreover, the coefficient of the highest-weight monomial \(x^\mu\) in the Schur polynomial \(s_\mu(x)\) is \(1\), so
\[
s_\mu(x)\ge x^\mu.
\]
Therefore
\begin{equation}\label{eq:tail_weight_bound}
    p_\mu(\delta)=\frac{x^{\mu-\delta}}{s_\mu(x)} \le p^{|\delta|}.
\end{equation}

Next define the total multiplicity at depth \(t\):
\begin{equation}
    D_\mu(t):= \sum_{|\delta|=t} m_\mu(\delta).
\end{equation}
Since the finite-dimensional irrep \(\h_\mu\) is a quotient of the Verma module \(M(\mu)\), each weight multiplicity is bounded above by Kostant's partition function:
\[
m_\mu(\delta)\le P_r(\delta),
\]
where \(P_r(\delta)\) counts the number of decompositions
\[
\delta=\sum_{\alpha\in \Phi_r^+} n_\alpha\,\alpha,
\qquad
n_\alpha\in \mathbb Z_{\ge 0},
\]
in terms of the positive roots of type \(A_{r-1}\).
Hence
\begin{equation}
    D_\mu(t)\le A_r(t),
    \qquad
    A_r(t):= \sum_{|\delta|=t} P_r(\delta).
\end{equation}

We now bound \(A_r(t)\) using its generating function. Since the height is additive,
\[
    \sum_{t\ge 0} A_r(t)\,q^t
    =
    \sum_{\delta\in Q_+} P_r(\delta)\, q^{|\delta|}
    =
    \prod_{\alpha\in \Phi_r^+}\frac{1}{1-q^{|\alpha|}}\ .
\]


For type \(A_{r-1}\) there are exactly \(r-h\) positive roots of height \(h\), so
\begin{equation}\label{eq:Ar_genfun}
    \sum_{t\ge 0} A_r(t)\,q^t
    =
    \prod_{h=1}^{r-1}\frac{1}{(1-q^h)^{\,r-h}}\ .
\end{equation}
Let
\begin{equation}
    N:= |\Phi_r^+|=\frac{r(r-1)}2\ .
\end{equation}
Because all coefficients are nonnegative, the generating function in~\eqref{eq:Ar_genfun} is coefficientwise bounded by
\[
    \prod_{h=1}^{r-1}\frac{1}{(1-q)^{\,r-h}}
    =
    \frac{1}{(1-q)^N}.
\]
Taking coefficients gives
\begin{equation}
    A_r(t)
    \le
    [q^t](1-q)^{-N}
    =
    \binom{t+N-1}{N-1}
    \le
    C_r\,(t+1)^{N-1},
\end{equation}
for some constant \(C_r>0\) depending only on \(r\). Thus
\begin{equation}\label{eq:Dt_bound}
    D_\mu(t)\le C_r\,(t+1)^\kappa,
    \qquad
    \kappa=N-1=\frac{r(r-1)}2-1.
\end{equation}

Combining~\eqref{eq:tail_weight_bound} and~\eqref{eq:Dt_bound}, we obtain
\begin{equation}
    \varepsilon_\mu(T)
    =
    \sum_{|\delta|>T} p_\mu(\delta)\,m_\mu(\delta)
    \le
    \sum_{t>T} D_\mu(t)\,p^t 
    \le
    C_r \sum_{t>T} (t+1)^\kappa p^t.
\end{equation}
Write \(t=T+1+s\) with \(s\ge 0\). Then
\[
    \sum_{t>T}(t+1)^\kappa p^t
    =
    p^{T+1}\sum_{s\ge 0}(T+s+2)^\kappa p^s
    \le
    p^{T+1}(T+2)^\kappa \sum_{s\ge 0}(s+2)^\kappa p^s.
\]
The series \(\sum_{s\ge 0}(s+2)^\kappa p^s\) converges because \(0<p<1\), so
\begin{equation}
    \varepsilon_\mu(T)=O\!\left(T^\kappa p^T\right).
\end{equation}
This proves the lemma.
\end{proof}





\begin{mlemma}{9}[Probability Ratio]
    Let \(\mu,\omega\) be two \(\mathrm{GL}(d,\mathbb C)\) irreps, and let \(x\) be a probability vector with \(r \le d\) strictly decreasing non-zero entries \(x_1>\cdots>x_r>0\). Let
    \[
    \xi := \min_{1 \le i \le r-1} (\ln x_i - \ln x_{i+1})
    \]
    be the logarithmic spectral gap. The ratio
    \begin{equation}
    Z(x):= \frac{x^\omega s_\mu(x)}{s_{\mu+\omega}(x)}
    \end{equation}
    satisfies
    \[
        1 \ge Z(x) \ge 1 - O\left(e^{-(g_\mu+1)\xi}\right)\ .
    \]
\end{mlemma} 

\begin{proof}
We first show the upper bound \(Z(x)\le 1\). Using the expansion of Schur polynomials into semistandard Young tableaux,
\begin{equation}
    s_\mu(x) = \sum_{u\prec\mu} K_{\mu,u}\,x^u \ .
\end{equation}
Multiplying by \(x^\omega\) shifts the weights:
\[
x^\omega s_\mu(x)=\sum_u K_{\mu,u}\,x^{u+\omega}\ .
\]
By Lemma~\ref{lem:kostka}, the Kostka numbers satisfy
\[
K_{\mu,u}\le K_{\mu+\omega,u+\omega}\ .
\]
Since the weights of \(\mu+\omega\) contain the shifted weights of \(\mu\) as a subset, every term in \(x^\omega s_\mu(x)\) appears in \(s_{\mu+\omega}(x)\) with coefficient at least as large. Thus
\[
s_{\mu+\omega}(x)-x^\omega s_\mu(x)\ge 0\ ,
\]
which implies \(Z(x)\le 1\).

For the lower bound, because \(x\) has exactly \(r\) non-zero entries, the Schur polynomials evaluate identically to those of \(\mathrm{GL}(r)\). Using the determinantal formula
\[
s_\lambda(x)=\frac{\det(x_i^{\lambda_j+r-j})}{\det(x_i^{r-j})}\ ,
\]
the denominators cancel in the ratio \(Z(x)\), leaving
\[
    Z(x) =
    x^\omega \frac{\det(x_i^{\ell_j})}{\det(x_i^{L_j})}\ ,
\]
where
\[
\ell_j=\mu_j+r-j,
\qquad
L_j=(\mu+\omega)_j+r-j=\ell_j+\omega_j\ .
\]
Expanding the determinants over the symmetric group \(S_r\),
\begin{equation}
    \det(x_i^{\ell_j})
    = \sum_{\sigma\in S_r}\mathrm{sgn}(\sigma)\prod_{j=1}^r x_j^{\ell_{\sigma(j)}}
    = \sum_{\sigma\in S_r}\mathrm{sgn}(\sigma)
    e^{\sum_j \ell_{\sigma(j)}\ln x_j}\ .
\end{equation}
Since both \(\ell\) and \(\ln x\) are strictly decreasing sequences, the dominant contribution comes from the identity permutation. For a single adjacent transposition \(\tau=(k,k+1)\), the exponent changes by
\[
    \Delta E
    = -(\ell_k-\ell_{k+1})(\ln x_k-\ln x_{k+1})\ .
\]
Now
\[
\ell_k-\ell_{k+1}=\mu_k-\mu_{k+1}+1\ge g_\mu+1\ ,
\]
so the gap is bounded by
\[
\epsilon_\mu:= (g_\mu+1)\xi\ .
\]
Therefore
\begin{equation}
    \det(x_i^{\ell_j})
    \ge
    x^\ell \left(1-(r!-1)e^{-\epsilon_\mu}\right)\ .
\end{equation}
Similarly, for the denominator,
\begin{equation}
    \det(x_i^{L_j})
    \le
    x^L \left(1+(r!-1)e^{-\epsilon_{\mu+\omega}}\right).
\end{equation}
Since \(\omega\) is dominant, \(g_{\mu+\omega}\ge g_\mu\), hence \(\epsilon_{\mu+\omega}\ge \epsilon_\mu\). Combining these estimates,
\[
    Z(x)
    \ge
    x^\omega \frac{x^\ell (1-(r!-1)e^{-\epsilon_\mu})}{x^L (1+(r!-1)e^{-\epsilon_\mu})}
    =
    \frac{1-(r!-1)e^{-\epsilon_\mu}}{1+(r!-1)e^{-\epsilon_\mu}}\ .
\]
Using \((1-y)/(1+y)\ge 1-2y\), we obtain
\[
    Z(x)\ge 1-O\left(e^{-(g_\mu+1)\xi}\right).
\]




\end{proof}



\begin{mlemma}{10}[Stability of the Weyl dimension]
Let \(\mu,\nu\) be dominant highest weights of \(\mathrm{GL}(d)\) with
\(\nu=\mu+\omega\), where \(\omega\) is dominant. Assume that \(\mu\) has at most \(r\) nonzero rows, and that
\[
g_\mu=\min_{1\le i\le r-1}(\mu_i-\mu_{i+1})=\Theta(n),
\qquad
\mu_r=\Theta(n),
\qquad
|\omega|=o(n).
\]
Then
\begin{equation}
    \frac{d_\mu}{d_\nu}
    = 1-O\!\left(\frac{|\omega|}{n}\right)
    = 1-O\!\left(\frac{d(\mu,\nu)}{n}\right).
\end{equation}
\end{mlemma}

\begin{proof}
By the Weyl dimension formula,
\begin{equation}
    \frac{d_\nu}{d_\mu}
    = \prod_{1\le i<j\le d}
    \frac{\mu_i-\mu_j+\omega_i-\omega_j+j-i}{\mu_i-\mu_j+j-i}
    = \prod_{1\le i<j\le d} \left(1+y_{ij}\right)\ ,
\end{equation}
where
\[
y_{ij}:=\frac{\omega_i-\omega_j}{\mu_i-\mu_j+j-i}\ .
\]
Since \(\omega\) is dominant, we have \(y_{ij}\ge 0\) for all \(i<j\).

We now bound the denominators from below. If \(1\le i<j\le r\), then
\[
\mu_i-\mu_j+j-i \ge \mu_i-\mu_j \ge g_\mu = \Theta(n)\ .
\]
If \(1\le i\le r<j\le d\), then \(\mu_j=0\) and
\[
\mu_i-\mu_j+j-i = \mu_i+j-i \ge \mu_r = \Theta(n).
\]
If \(r<i<j\le d\), then \(\mu_i=\mu_j=0\) and also \(\omega_i=\omega_j=0\), so
\(y_{ij}=0\). Therefore
\begin{equation}
    0\le y_{ij}\le C_1\frac{|\omega|}{n}
\end{equation}
for all \(i<j\), where \(C_1\) depends only on \(d\).

Since \(d\) is fixed, there are only \(O_d(1)\) factors, and
\begin{equation}
    \sum_{i<j} y_{ij}
    \le \frac{C_2}{n}\sum_{i<j}(\omega_i-\omega_j)
    \le C_3\frac{|\omega|}{n}.
\end{equation}
Because \(|\omega|=o(n)\), we have \(y_{ij}=o(1)\) uniformly. Hence
\[
\log(1+y_{ij})=y_{ij}+O(y_{ij}^2),
\]
and summing over \(i<j\) gives
\begin{equation}
    \log\frac{d_\nu}{d_\mu}
    = \sum_{i<j}y_{ij}
    + O\!\left(\sum_{i<j}y_{ij}^2\right)
    = \sum_{1\le i<j\le d}
    \frac{\omega_i-\omega_j}{\mu_i-\mu_j+j-i}
    + O\!\left(\frac{|\omega|^2}{n^2}\right).
\end{equation}
Exponentiating yields
\[
\frac{d_\nu}{d_\mu}=1+O\!\left(\frac{|\omega|}{n}\right).
\]
\end{proof}

\bibliographystyle{utphys}
\bibliography{free}

\end{document}